\begin{ex}%[1D5V1-2]
% nguon: Archive/SP_SauLuoi_2024-2025/SP_HKI/1-TN-DS-TLN-THPT-NguyenThiMinhKhai-HCM-HKI-NH24-25.tex #8
%[1-TN-DS-TLN-THPT-NguyenThiMinhKhai-HCM-HKI-NH24-25]%[Đức Dương]
	Bảng thống kê doanh số bán hàng của các nhân viên một siêu thị điện máy trong một dịp Black Friday như sau:
	\begin{center}
		\begin{tabular}{|c|c|c|c|c|c|}
			\hline \begin{tabular}{c} 
				Doanh số \\
				(triệu đồng)
			\end{tabular} & $[20;30)$ & $[30;40)$ & $[40;50)$ & $[50;60)$ & $[60;70)$ \\
			\hline Số nhân viên & $4$ & $8$ & $12$ & $7$ & $5$ \\
			\hline
		\end{tabular}
	\end{center}
	\choiceTF
	{Cỡ mẫu của mẫu số liệu trên là $n=35$}
	{\True Số trung bình của mẫu số liệu trên là $\overline{x}\approx 45{,}28$ triệu đồng (\textit{kết quả làm tròn đến hàng phần trăm})}
	{Tứ phân vị thứ nhất của mẫu số liệu trên là $Q_1=36{,}5$}
	{Cửa hàng dự định sẽ thưởng $25\%$ số nhân viên có doanh số bán hàng cao nhất. Theo mẫu sổ liệu trên, của hàng nên khen thưởng các nhân viên có doanh số bán hàng ít nhất từ $53{,}29$ triệu đồng (\textit{kết quả làm tròn đến hàng phần trăm})}
	\loigiai{
		Theo bài ra ta có bảng tần số ghép nhóm sau
		\begin{center}
			\begin{tabular}{|c|c|c|c|c|c|}
				\hline \begin{tabular}{c} 
					Doanh số \\
					(triệu đồng)
				\end{tabular} & $[20;30)$ & $[30;40)$ & $[40;50)$ & $[50;60)$ & $[60;70)$ \\
				\hline Giá trị đại diện & $25$ & $35$ & $45$ & $55$ & $65$ \\		
				\hline Số nhân viên & $4$ & $8$ & $12$ & $7$ & $5$ \\
				\hline
			\end{tabular}
		\end{center}
		\begin{itemchoice}
			\itemch \textbf{Sai}.\\
			Cỡ mẫu của mẫu số liệu trên là $n=4+8+12+7+5=36$.
			\itemch \textbf{Đúng}.\\
			Giá trị trung bình của mẫu số liệu 
			\[ \overline{x}=\dfrac{25\cdot 4+35\cdot 8+45\cdot 12+55\cdot 7+65\cdot 5}{36}=\dfrac{815}{18}\approx 45{,}28.\]
			\itemch \textbf{Sai}.\\
			Ta có $\dfrac{n}{4}=9$ mà $4<9<12$ nên suy ra tứ phân vị thứ nhất là 
			\[ Q_1=30+\dfrac{9-4}{8}\cdot 10=36{,}25.\]
			\itemch \textbf{Sai}.\\
			Ta có $\dfrac{3n}{4}=27$ mà $4+8+12<27<4+8+12+7$ nên suy ra tứ phân vị thứ ba là 
			\[ Q_3=50+\dfrac{27-24}{7}\cdot 10=\dfrac{380}{7}\approx 54{,}29. \]
			Vậy cửa hàng nên khen thưởng các nhân viên có doanh số bán hàng ít nhất là $54{,}29$ triệu đồng.
		\end{itemchoice}
	}
\end{ex}

\begin{ex}%[1D5H2-2]
% nguon: Archive/SP_SauLuoi_2024-2025/SP_HKI/1-TN-DS-TLN-THPT-ChuyenHungVuong-PhuTho-HKI-NH24-25.tex #15
%[11 - Kiểm tra học kì 1-2024-2025-Chuyên Hùng Vương-Phú Thọ]%[Trần Văn Hùng]
	Cho mẫu số liệu ghép nhóm về Điểm khảo sát chất lượng lần $1$ môn Toán của một lớp $11$ trường THPT chuyên Hùng Vương năm học $2024-2025$ như sau
	\begin{center}
		\begin{tabular}{|c|c|c|c|c|c|c|}
			\hline
			Điểm  &$[4;5)$  &$[5;6)$&$[6;7)$  &$[7;8)$  &$[8;9)$  &$[9;10)$  \\
			\hline
			Số học sinh  &$1$  &$6$ &$11$ &$15$  &$9$  &$3$  \\
			\hline
		\end{tabular}
	\end{center}
	\choiceTF
	{Cỡ mẫu là $44$}
	{\True Mốt và số trung vị của mẫu số liệu thuộc cùng một nhóm}
	{Số trung bình của mẫu số liệu là $7{,}2$}
	{\True Có khoảng $50$\% học sinh trong lớp đạt từ $7{,}3$ điểm trở lên}
	\loigiai{
		\begin{center}
			\begin{tabular}{|c|c|c|c|c|c|c|}
				\hline
				Điểm  &$[4;5)$  &$[5;6)$&$[6;7)$  &$[7;8)$  &$[8;9)$  &$[9;10)$  \\
				\hline
				Giá trị đại diện &$4{,}5$  &$5{,}5$&$6{,}5$  &$7{,}5$  &$8{,}5$  &$9{,}5$  \\
				\hline
				Số học sinh  &$1$  &$6$ &$11$ &$15$  &$9$  &$3$  \\
				\hline
				Tần số tích lũy  &$1$  &$7$ &$18$ &$33$  &$42$  &$45$  \\
				\hline
			\end{tabular}
		\end{center}
		\begin{itemchoice}
			\itemch \textbf{Sai}.\\
			Cỡ mẫu là $n=1+6+11+15+9+3=45$.
			\itemch \textbf{Đúng}.\\
			Nhóm có tần số lớn nhất là $[7,8)$, suy ra mốt thuộc nhóm $[7;8)$.\\
			Vì $18<\dfrac{n}{2}=\dfrac{45}{2}<33$ nên 
			$M_e=6+\dfrac{\dfrac{45}{2}-18}{15}\cdot (7-6)=\dfrac{73}{10}=7{,}3$.\\
			Vậy mốt và số trung vị của mẫu số liệu thuộc cùng một nhóm $[7;8)$.
			\itemch \textbf{Sai}.\\
			Số trung bình $\overline{x}=\dfrac{1\cdot 4{,}5+6\cdot 5{,}5+11\cdot 6{,}5+15\cdot 7{,}5+9\cdot 8{,}5+3\cdot 9{,}5}{45}=\dfrac{653}{90}\approx 7{,}3$.
			\itemch \textbf{Đúng}.\\
			Vì $M_e=7{,}3$ nên có khoảng $50$\% học sinh trong lớp đạt từ $7{,}3$ điểm trở lên.
		\end{itemchoice}
	}
\end{ex}

\begin{bt}%[1D5H1-4]
% nguon: Archive/SP_SauLuoi_2024-2025/SP_HKI/1-TN-DS-TLN-THPT-ChuyenHungVuong-PhuTho-HKI-NH24-25.tex #20
%[Dự án đề kiểm tra Toán khối 11 học kì 1 NH24-25-Dot 5-Thái Bảo ]%[THPT Chuyên Hùng Vương - Phú Thọ]
	Người ta đo chiều cao của $35$ cây bạch đàn (đơn vị: m) và thu được kết quả như sau
	\begin{center}
		\begin{tabular}{|l|l|l|l|l|l|l|l|l|l|}
			\hline
			$6{,}6$ & $6{,}9$ & $7{,}0$ & $7{,}2$ & $7{,}3$ & $7{,}4$ & $7{,}5$ & $7{,}5$ & $7{,}6$ & $7{,}7$ \\
			\hline
			$7{,}7$ & $7{,}8$ & $7{,}8$ & $7{,}9$ & $7{,}9$ & $8{,}0$ & $8{,}0$ & $8{,}0$ & $8{,}1$ & $8{,}1$ \\
			\hline
			$8{,}2$ & $8{,}2$ & $8{,}2$ & $8{,}3$ & $8{,}3$ & $8{,}4$ & $8{,}5$ & $8{,}5$ & $8{,}6$ & $8{,}7$ \\
			\hline
			$8{,}7$ & $8{,}8$ & $8{,}9$ & $9{,}0$ & $9{,}4$ & & & & & \\
			\hline
		\end{tabular}
	\end{center}
	Từ mẫu số liệu không ghép nhóm trên, hãy ghép các số liệu thành $6$ nhóm có độ dài bằng nhau trong đó có nhóm $[6{,}5; 7{,}0)$. Uớc lượng mốt của mẫu số liệu ghép nhóm (kết quả làm tròn đến hàng phần trăm).
	\shortans{$8{,}17$}
	\loigiai{
		Ghép các số liệu thành $6$ nhóm có độ dài bằng nhau trong đó có nhóm $[6{,}5; 7{,}0)$ ta được mẫu số liệu ghép nhóm sau
		\begin{center}
			\begin{tabular}{|c|c|c|c|c|c|c|}
				\hline
				Chiều cao (m) & $[6{,}5;7{,}0)$ & $[7{,}0;7{,}5)$ & $[7{,}5;8{,}0)$ & $[8{,}0;8{,}5)$ & $[8{,}5;9{,}0)$ & $[9{,}0;9{,}5)$ \\
				\hline
				Số cây & $2$ & $4$ & $9$ & $11$ & $7$ & $2$ \\
				\hline
			\end{tabular}
		\end{center}
		Tần số lớn nhất của mẫu dữ liệu ghép nhóm trên là $11$ nên nhóm chứa mốt là nhóm $[8{,}0;8{,}5)$. Ta có $j=4$, $a_4=8{,}0$, $m_4=11$, $m_3=9$, $m_5=7$, $h=0{,}5$. Do đó
		$$M_o=a_j+\dfrac{m_j-m_{j-1}}{\left(m_j-m_{j-1}\right)+\left(m_j-m_{j+1}\right)} \cdot h=8+\dfrac{11-9}{(11-9)+(11-7)}\cdot 0{,}5=\dfrac{49}{6}\approx8{,}17.$$
		
	}
\end{bt}

\begin{ex}%[1D5V2-3]
% nguon: Archive/SP_SauLuoi_2024-2025/SP_HKI/1-TN-DS-TLN-THPT-LeQuyDon-HCM-HKI-NH24-25.tex #28
%[Dự án đề kiểm tra Toán khối 11 HKI NH24-25-Dot 5- Sy Truong]%[THPT Lê Quý Đôn-TP HCM]
	Người ta đo đường kính của một số cây gỗ được trồng sau $ 12 $ năm (đơn vị: centimét), họ thu được bảng tần số ghép nhóm sau:
	\begin{center}
		\begin{tabular}{|c|c|c|c|c|c|}
			\hline Đường kính (cm) &{$[20; 25)$} &{$[25; 30)$} &{$[30; 35)$} &{$[35; 40)$} &{$[40; 45)$} \\
			\hline Số cây & $ 4 $ & $ 12 $ & $ 26 $ & $ 13 $ & $ 6 $ \\
			\hline
		\end{tabular}
	\end{center}
	\choiceTF
	{\True Cỡ mẫu của mẫu số liệu là $n=61$}
	{\True Tứ phân vị thứ ba của mẫu số liệu ghép nhóm là $Q_3 \approx 36,44$}
	{\True Tứ phân vị thứ hai của mẫu số liệu ghép nhóm là $Q_2 \approx 32,79$}
	{Tứ phân vị thứ nhất của mẫu số liệu ghép nhóm là $Q_1=28,69$}
	\loigiai{
		\begin{itemchoice}
			\itemch \textbf{Đúng}.
			Ta có cỡ mẫu là $ n= 4 + 12 + 26 + 13 + 6 = 61$.
			\itemch \textbf{Đúng}.
			Do $ n=61 $, nên tứ phân vị thứ ba là $ \dfrac{1}{2}\left(x_{45}+x_{46}\right) $.\\
			Mà $ x_{45},x_{46}\in [35;40) $ nên ta có
			$$Q_3=35+\dfrac{\tfrac{3\cdot61}{4}-(4+12+26)}{13}\left(40-35\right)\approx 36{,}44.$$
			\itemch \textbf{Đúng}.
			Do $ n=61 $ nên trung vị là giá trị $x_{31}\in [30;35)$, do đó
			$$Q_2=30+\dfrac{\tfrac{61}{2}-(4+12)}{26}\left(40-35\right)\approx 32{,}79.$$
			\itemch \textbf{Sai}.
			Do $n=61$, nên tứ phân vị thứ nhất là $\dfrac{1}{2}\left(x_{15}+x_{16}\right)$.\\
			Mà $ x_{15}$, $x_{16}\in [25;30) $ nên ta có
			$$Q_1=25+\dfrac{\tfrac{61}{4}-4}{12}\left(30-25\right)\approx 29{,69}.$$
		\end{itemchoice}
	}
\end{ex}

\begin{bt}%[1D5H2-2]
% nguon: Archive/SP_SauLuoi_2024-2025/SP_HKI/1-TN-DS-TLN-THPT-LeQuyDon-HCM-HKI-NH24-25.tex #30
%[Dự án đề kiểm tra Toán 11 HKI NH24-25-Võ Thị Thùy Trang]%[Lê Quý Đôn - Tp HCM]
	Trong một hội thao, thời gian chạy $200$ m của một nhóm các vận động viên được ghi lại ở bảng sau
	\begin{center}
		\begin{tabular}{|c|c|c|c|c|c|}
			\hline
			Thời gian (giây) &$[21;21{,}5)$& $[21{,}5;22)$ &$[22;22{,}5)$&$[22{,}5;23)$ &$[23;23{,}5)$\\
			\hline
			Số vận động viên & $5$ & $12$ & $32$ & $45$ & $30$ \\
			\hline
		\end{tabular}
	\end{center}
	Dựa vào bảng số liệu trên, ban tổ chức muốn chọn ra khoảng $50\%$ số vận động viên chạy nhanh nhất để tiếp tục thi vòng $2$. Ban tổ chức nên chọn các vận động viên có thời gian chạy không quá bao nhiêu giây (kết quả làm tròn đến hàng phần chục)?
	\shortans{$2{,}64$}
	\loigiai{Trung vị của bảng số liệu trên là $M_\text{e}=u_m+\dfrac{\dfrac{n}{2}-C}{n_m}\left(u_{m+1}-u_m\right)$.\\
		$M_\text{e}=22{,}5+\dfrac{\dfrac{124}{2}-49}{45}\left(23-22{,}5\right)=\dfrac{1019}{45}\approx 22{,}64$ (giây).}
\end{bt}

\begin{bt}%[1D5H2-3]
% nguon: Archive/SP_SauLuoi_2024-2025/SP_HKI/2-TN-DS-TLN-THPT-ThiXaQuangTri-QuangTri-HKI-NH24-25.tex #36
%[Dự án đề kiểm tra Toán 12 HKI NH23-24- Xuan Vy Pham]%[THPT Thị Xã Quảng Trị]
Mẫu số liệu ghép nhóm thống kê mức lương của một công ty (đơn vị: triệu đồng) được cho trong bảng dưới đây.
\begin{center}
	\begin{tabular}{|c|c|c|c|c|c|c|}
\hline
Nhóm&$[6;8)$&$[8;10)$&$[10;12)$&$[12;14)$&$[14;16)$ & \\
\hline
Tần số& $6$&$14$&$18$&$10$&$2$&$n=50$ \\
\hline
\end{tabular}
\end{center}
Tìm tứ phân vị thứ nhất của mẫu số liệu ghép nhóm \textit{(làm tròn đến hàng phần trăm)}.
\shortans{8,93}
\loigiai{
Gọi $x_1$, $x_2$,$\ldots$, $x_{50}$ là các mức lương của một công ty được sắp xếp theo thứ tự không giảm.\\
Ta có $x_1$, $x_2$,$\ldots$, $x_6 \in [6;8)$, $x_7,\ldots,x_{20} \in [8;10)$, $x_{21},\ldots,x_{38} \in [10;12)$, $x_{39},\ldots,x_{48} \in [12;14)$, $x_{49}$, $x_{50} \in [14;16)$.\\
Tứ phân vị thứ nhất là $x_{13}$ nên nhóm chứa tứ phân vị thứ nhất là $[8;10)$.\\
Do đó $Q_1=8+\dfrac{\dfrac{50}{4}-6}{14} \cdot (10-8)=\dfrac{125}{14}\approx 8{,}93$.
}
\end{bt}

\begin{ex}%[1D5V2-3]
% nguon: Archive/SP_SauLuoi_2024-2025/SP_HKI/1-TL-TN-DS-THPT-Hoang-Viet-Daklak-HKI-NH24-25.tex #44
%[TH-THCS-THPT-HoangViet-DakLak-HKI-NH24-25]%[Trần Hưng]
	Một bảng xếp hạng đã tính điểm chuẩn hoá cho chỉ số nghiên cứu của một số trường đại học ở Việt Nam và thu được kết quả sau:
	\begin{center}
		\begin{tabular}{|c|c|c|c|c|c|c|}
			\hline
			Điểm & $[10;20)$ & $[20;30)$ & $[30;40)$ & $[40;50)$ & $[50,60)$ & $[60,70)$\\
			\hline
			Số trường & $4$ & $19$ & $6$ & $2$ & $3$ & $1$\\
			\hline
		\end{tabular}
	\end{center}
	\choiceTF
	{\True Số liệu đã cho cho có $35$ mẫu số liệu}
	{Số trung bình của mẫu số liệu đã cho là $28$}
	{Số trung vị của mẫu số liệu là $M_{e}=12$}
	{\True Ngưỡng điểm để đưa ra danh sách $25\%$ trường đại học có chỉ số nghiên cứu tốt nhất Việt Nam là trên $35{,}42$}
	\loigiai{
		\begin{itemchoice}
			\itemch 
			Số các số liệu đã cho có mẫu số liệu là $n=4+19+6+2+3+1=35$.
			\itemch 
			Số trung bình cộng của mẫu số liệu đã cho là\\
			\[\overline{x}=\dfrac{4\cdot 15+19\cdot 25+6\cdot 35+2\cdot 45+3\cdot 55+1\cdot 65}{35}\approx30{,}43.\]
			\itemch 
			Vị trí của số trung vị là $\dfrac{35+1}{2}=18$.\\
			Do đó, giá trị trung vị của mẫu số liệu nằm trong khoảng $[20;30)$.
			\itemch 
			Gọi $x_1$, $x_2$, $\ldots$, $x_{35}$ là điểm chuẩn hóa cho chỉ số nghiên cứu của các trường đại học và giả sử dãy này đã sắp xếp theo thứ tự tăng dần. Khi đó, trung vị của mẫu số liệu là $M_e=x_{18}$ và tứ phân vị thứ ba $Q_3$ của mẫu số liệu là trung vị của nửa số liệu bên phải $M_e$, đó là dãy gồm 17 số liệu $x_{19}$, $x_{20}$, $\ldots$, $x_{35}$, do đó $Q_3=x_{27}$.\\
			Do $x_{27}$ thuộc nhóm $[30;40$ ) nên nhóm này chứa $Q_3$.\\
			Do đó, $p=3$; $a_3=30$; $m_3=6$; $m_1+m_2=4+19=23$; $a_4-a_3=40-30=10$ và ta có 		
			\[Q_3=30+\dfrac{\dfrac{3{,}35}{4}-23}{6} \cdot 10 \approx 35{,}42.\]	
			$25\%$ trường đại học có chỉ số nghiên cứu tốt nhất là $25\%$ giá trị nằm bên phải $Q_3$.\\
			Vậy điểm ngưỡng để đưa ra danh sách $25\%$ trường đại học có chỉ số nghiên cứu tốt nhất Việt Nam là những trường có điểm chuẩn hóa lớn hơn hoặc bằng $35{,}42$.
		\end{itemchoice}
	}
\end{ex}

\begin{ex}%[1D5H2-2]
% nguon: Archive/SP_SauLuoi_2024-2025/SP_HKI/1-TN-DS-TLN-THPT-PhanBoiChau-HKI-BinhThuan-NH23-24.tex #61
%[Dự án đề kiểm tra Toán 11 HKI NH23-24- Xuan Vy Pham]%[THPT Phan Bội Châu]
Theo thống kê điểm thi trung bình môn toán của một nhóm học sinh đã trúng tuyển vào THPT A trong kỳ thi tuyển sinh 10 cho kết quả như bảng sau
\begin{longtable}{|c|c|c|c|c|c|c|c|}
	\hline
	Khoảng điểm&$[6{,}5;7)$&$[7;7{,}5)$&$[7{,}5;8)$&$[8;8{,}5)$&$[8{,}5;9)$ &$[9;9{,}5)$&$[9{,}5;10)$  \\
	\hline
	Tần số& $7$&$10$&$17$&$24$&$13$&$8$&$5$  \\
	\hline
\end{longtable}
\choiceTF
{Số học sinh được khảo sát là $86$}
{Điểm trung bình môn Toán của nhóm học sinh đã trúng tuyển vào trường THPT A trong kì thi tuyển sinh vào lớp 10 nhỏ hơn $7$}
{\True Mốt của mẫu số liệu ghép nhóm trên là $\dfrac{295}{36}$}
{\True Trung vị của mẫu số liệu ghép nhóm trên là $\dfrac{49}{6}$}
\loigiai{
\begin{itemchoice}
	\itemch \textbf{Sai.} Vì cỡ mẫu là $n=7+10+17+24+13+8+5=84$ nên số học sinh được khảo sát là $84$.
	\itemch	\textbf{Sai.} Bảng tần số ghép nhóm thêm hàng giá trị đại diện có dạng
	\begin{longtable}{|c|c|c|c|c|c|c|c|}
		\hline
		Khoảng điểm&$[6{,}5;7)$&$[7;7{,}5)$&$[7{,}5;8)$&$[8;8{,}5)$&$[8{,}5;9)$ &$[9;9{,}5)$&$[9{,}5;10)$  \\
		\hline
		Giá trị đại diện&$6{,}75$&$7{,}25$&$7{,}75$&$8{,}25$&$8{,}75$&$9{,}25$&$9{,}75$\\
		\hline
		Tần số&$7$&$10$&$17$&$24$&$13$&$8$&$5$  \\
		\hline
	\end{longtable}
	Điểm trung bình môn Toán của nhóm học sinh đã trúng tuyển vào trường THPT A trong kì thi tuyển sinh vào lớp 10 là
	$$\overline{x}=\dfrac{6{,}75 \cdot 7+ 7{,}25 \cdot 10+ 7{,}75 \cdot 17+ 8{,}25 \cdot 24+8{,}75 \cdot 13+ 9{,}25 \cdot 8 + 9{,}75\cdot 5}{84}=\dfrac{49}{6}\approx 8{,}17>7.$$ 
	\itemch	\textbf{Đúng.} Nhóm chứa mốt là $[8;8{,}5)$. Vậy mốt của mẫu số liệu ghép nhóm trên là 
	$$M_o=8+\dfrac{24-17}{(24-17)+(24-13)} \cdot (8{,}5-8)=\dfrac{295}{36}.$$
	\itemch \textbf{Đúng.} Nhóm chứa trung vị là $[8;8{,}5)$. Do đó trung vị của mẫu số liệu ghép nhóm trên là 
	$$M_e=8+\dfrac{\dfrac{84}{2}-(7+10+17)}{24} \cdot (8{,}5-8)=\dfrac{49}{6}.$$
\end{itemchoice}
}
\end{ex}

\begin{bt}%[1D5H1-3]
% nguon: Archive/SP_SauLuoi_2024-2025/SP_HKI/1-TN-DS-TLN-THPT-PhanBoiChau-HKI-BinhThuan-NH23-24.tex #66
%[KNTT - Lớp 12 - Ôn tập Cuối kì 1 - Đề Phan Bội Châu-Bình Thuận]%[Phạm Văn Phú]
    Lớp 12A có 50 học sinh, mỗi học sinh làm bài thi trắc nghiệm môn Toán có 50 câu. Biết số câu trả lời đúng bài thi trắc nghiệm môn Toán được cho bởi bảng sau
    \begin{center}
        \begin{tabular}{|l|c|c|c|c|c|}
            \hline
            Số câu đúng & $[14;21)$ & $[21;28)$ & $[28;35)$ & $[35;42)$ & $[42;49)$ \\
            \hline
            Số học sinh & $4$ & $x$ & $25$ & $y$ & $7$ \\
            \hline
        \end{tabular}
    \end{center}
    Tính $3x-2y$ biết số trung bình của mẫu số liệu bằng $32,06$
    \shortans{12}
    \loigiai{
        Tổng số học sinh là $x+y+36=50$, suy ra $x+y=14$. \,(1)\\
        Lại có $\dfrac{ 4.17,5+x.24,5+25.31,5+y.38,5+7.45,5}{50}=32,06$.\\
        Từ đó suy ra $24,5x+38,5y=427$. \,(2)\\
        Từ (1) và (2) ta được $x=8$ và $y=6$ suy ra $3x-2y=12$
    }
\end{bt}

\begin{bt}%[1D5H2-3]
% nguon: Archive/SP_SauLuoi_2024-2025/SP_HKI/1-TN-DS-TLN-THPT-ChuyenLeHongPhong-TpHCM-HKI-NH24-25.tex #85
[Dự án đề kiểm tra môn Toán khối 11 HKI NH24-25 - Đợt 5 - Hieu Hieu Minh Minh]%[THPT Chuyên Lê Hồng Phong - Tp HCM]
	Thời gian (phút) truy cập Internet mỗi buổi tối của một số học sinh được cho trong bảng sau:
	\begin{longtable}{|c|c|c|c|c|c|}
		\hline
		Thời gian (giờ)&$[9{,}5; 12{,}5)$ & $[12{,}5; 15{,}5)$ & $[15{,}5;18{,}5)$ & $[18{,}5;21{,}5)$ &$[21{,}5;24{,}5)$  \\
		\hline
		Số học sinh& $3$ & $12$ & $15$ & $24$ & $2$  \\
		\hline
	\end{longtable}
	Tính tứ phân vị thứ nhất của mẫu số liệu ghép nhóm ở trên (làm tròn kết quả tới hàng phần mười).
	\loigiai{
		Cỡ mẫu là $n=56$.\\
		Tứ phân vị thứ nhất $Q_1$ là $\dfrac{x_{14}+x_{15}}{2}$. \\
		Do $x_{14}$, $x_{15}$ đều thuộc nhóm $[12{,}5 ; 15{,}5)$ nên nhóm này chứa $Q_1$.\\
		Do đó $Q_1=12{,}5+\dfrac{\frac{56}{4}-3}{12} \cdot 3\approx 15{,}3$.
	}
\end{bt}

\begin{bt}%[1D5H1-4]
% nguon: Archive/SP_SauLuoi_2024-2025/Sp_GHKII/data/GHKII/1-TN-DS-TLN-TL-THPT-HoaHoi-BRVT-GHKII-NH24-25.tex #159
%[Dự án tex đề GHKII THPT Hoà Hội - Bà Rịa Vũng Tàu- 24-25][Đỗ Minh Vũ]
		 		Bảng số liệu ghép nhóm sau cho biết chiều cao (cm) của $50$ học sinh lớp $11$A $\colon$		 		
		 		\begin{center}
		 			\begin{tabular}{|c|c|c|c|c|c|}
		 				\hline
		 				Khoảng chiều cao (cm) & $[145;150)$ & $[150;155)$ & $[155;160)$ & $[160;165)$ & $[165;170)$\\
		 				\hline
		 				Số học sinh & $7$ & $14$ & $10$ & $10$ & $9$ \\
		 				\hline
		 			\end{tabular}
		 		\end{center}
		 		Tính mốt của mẫu số liệu trên (\textit{kết quả làm tròn đến hàng đơn vị}).
		 		\shortans[oly]{$153$}
		 	\loigiai{
		 			Để tìm mốt của mẫu số liệu ghép nhóm, ta thực hiện theo các bước sau
		 		\begin{itemize}
		 			\item Bước 1. Xác định nhóm có tần số lớn nhất (gọi là nhóm chứa mốt), ở đây là giả sử là nhóm $j\colon$ $\left[a_j;a_{j+1}\right)$.
		 			\item Bước 2. Mốt được xác định là
		 			$$M_0=a_j+\frac{m_j-m_{j-1}}{\left(m_j-m_{j-1}\right)+\left(m_j-m_{j+1}\right)}\cdot h$$ trong đó $m_j$ là tần số của nhóm $j$ (quy ước $m_0=m_{k+1}=0$ ) và $h$ là độ dài của nhóm.
		 		\end{itemize}
		 	\textbf{Áp dụng}
		 		\begin{itemize}
		 			\item Nhóm chứa mốt là $[150;155)$.
		 			\item $M_0=150+\dfrac{14-7}{(14-7)+(14-10)}\cdot (155-150)=\dfrac{1685}{11} \approx 153$ (cm).
		 			
		 		\end{itemize}
		 	}
		 \end{bt}

\begin{ex}%[1D5N2-3]
% nguon: Archive/SP_SauLuoi_2024-2025/Sp_GHKII/data/GHKII/1-TL-TN-DS-THPT-Tran-Quang-Khai-GHKII-NH24-25.tex #183
%[Dự án đề kiểm tra Toán 11 GHKII NH24-25- Hải Phụng]%[THPT Trần Quang Khải - BRVT]
	Khảo sát thời gian tập thể dục của một số học sinh khối 12 thu được mẫu số liệu ghép nhóm sau:
	\begin{center}
		\begin{tabular}{|c|c|c|c|c|c|}
			\hline
			Thời gian & $[0;20)$ & $[20;40)$ & $[40;60)$ & $[60;80)$ & $[80;100)$ \\
			\hline
			Số học sinh & 5 & 9 & 12 & 10 & 6 \\
			\hline
		\end{tabular}
	\end{center}
	\choiceTF
	{Nhóm $[20;40)$ chứa tứ phân vị thứ hai}
	{\True Trung vị của mẫu số liệu trên bằng $\dfrac{155}{3}$}
	{\True Mốt của mẫu số liệu trên bằng $52$}
	{\True Đầu mút phải của nhóm $[40;60)$ bằng $60$}
	\loigiai{
		\begin{itemchoice}
			\itemch Ta có $n=5+9+12+10+6=42$ và $\dfrac{42+1}{2} = 21{,}5$ nên $Q_2\in [40;60)$.
			\itemch Trung vị là $Q_2=40+\dfrac{\frac{42}{2}-14}{12}\cdot 20 = \dfrac{155}{3}$.
			\itemch Vì $M_0 \in [40;60)$ nên $M_0 = 40 + \dfrac{12-9}{(12-9)+(12-10)}\cdot 20 = 52$.
			\itemch Đầu mút phải của nhóm $[40;60)$ bằng $60$.
		\end{itemchoice}
	}
\end{ex}

\begin{bt}%[1D6N3-5]
% nguon: Archive/SP_SauLuoi_2024-2025/Sp_GHKII/data/GHKII/1-TL-TN-DS-THPT-Tran-Quang-Khai-GHKII-NH24-25.tex #185
%[Dự án đề kiểm tra Toán 11 GHKII NH24-25- Phạm Phú]%[THPT Trần Quang Khải - BRVT]
	Trong khoa học môi trường người ta sử dụng công thức $A=k\cdot D^b$ để ước tính tuổi của một cây dựa vào đường kính của thân cây, trong đó: $A$ là tuổi của cây (tính bằng năm), $D$ là đường kính thân cây (tính bằng centimet), $k$ và $b$ là các hằng số phụ thuộc vào loại cây. Với một loại cây rừng đường kính $30$ cm, các nhà nghiên cứu xác định được rằng $k=2,5$ và $b=1,3$. Hãy tính tuổi của cây rừng trên. {\it (làm tròn kết quả đến hàng đơn vị)}.
	\shortans[]{$208$}
	\loigiai{
 	Từ công thức $A=k\cdot D^h$ ta có tuổi của cây cần tính là $A=2,5\cdot 30^{1,3}\approx 208$ tuổi.
	}
\end{bt}

\begin{bt}%[1D5N1-3]
% nguon: Archive/SP_SauLuoi_2024-2025/Sp_GHKII/data/GHKII/1-TL-TN-DS-THPT-Tran-Quang-Khai-GHKII-NH24-25.tex #187
%[Dự án đề kiểm tra Toán 11 GHKII NH24-25- Phạm Phú]%[THPT Trần Quang Khải - BRVT]
	 Cho mẫu số liệu ghép nhóm về lương (đơn vị: triệu đồng) và số nhân viên như sau
	\begin{center}
		\begin{tabular}{|c|c|c|c|c|c|c|}
			\hline Cân nặng ($m$) &{$[10;15)$}&{$[15;20)$}&{$[20;25)$}&{$[25;30)$}&{$[30;35)$}&{$[35;40)$}\\
			\hline Số nhân viên & 5 & 8 & 12 & 10 & 2 & 3 \\
			\hline
		\end{tabular}
	\end{center}
	Tính lương (đơn vị: triệu đồng) trung bình từ mẫu số liệu ghép nhóm trên {\it (kết quả làm tròn đến hàng phần chục)}.
	 \shortans[]{$23,1$}
	\loigiai{
	Lương trung bình của mẫu số liệu trên là $$\overline{x}=\dfrac{5\cdot 12,5+8\cdot 17,5+12\cdot 22,5+10\cdot 27,5+2\cdot 32,5+3\cdot 37,5}{40}\approx 23{,}1$$.
	}
\end{bt}

\begin{bt}%[1D5H1-3]
% nguon: Archive/SP_SauLuoi_2024-2025/Sp_GHKII/data/GHKII/1-TN-DS-TLNTHPT-Kon-Tum-GHKII-NH24-25.tex #202
%[1-HK1- THPT Kon-Tum - Năm học 2024-2025]%[Nguyễn Hoàng Việt - Đợt 7]
	Trường THPT Kon Tum đã tiến hành khảo sát mức độ hài lòng của học sinh đối với việc giảng dạy. Dữ liệu thu thập được như sau:
	\begin{center}
		\begin{tabular}{|c|c|c|c|c|c|}
			\hline
			Số điểm &  [50;60) & [60;70) & [70;80) & [80;90) & [90;100]  \\
			\hline
			Số người & 4 & 7 & 10 & 17 & 12 \\
			\hline
		\end{tabular}
	\end{center}
	Tính giá trị trung bình của mẫu số liệu ghép nhóm ở trên.
	\shortans{$80{,}2$}
	\loigiai{
		Giá trị trung bình $\overline{x} = \dfrac{55 \cdot 4 + 65 \cdot 7 + 75 \cdot 10 + 85 \cdot 17 + 95 \cdot 12}{50} = 80{,}2$.
	}
\end{bt}

\begin{ex}%[1D5H2-2]
% nguon: Archive/SP_SauLuoi_2024-2025/Sp_GHKII/data/GHKII/1-TN-DS-TLN-TL-THPT-Edison-HaiPhong-GKII-NH24-25.tex #254
%[Dự án tex đề GHKII THPT Edison - Hải Phòng- 24-25][Nguyễn Hữu Đức]
	Khi đo mắt cho học sinh khối $10$ ở một trường THPT, nhân viên y tế ghi nhận số liệu theo bảng sau
	\begin{center}
		\begin{tabular}{|c|c|c|c|c|c|}
			\hline \textbf{Thời gian đo} &$[0{,}25; 0{,}75)$&$[0{,}75; 1{,}25)$&$[1{,}25;1{,}75)$&$[1{,}75; 2{,}25)$&$[2{,}25; 2{,}75)$\\
			\hline \textbf{Số lần} &$25$ & $32$&$14$ &$12$ &$4$ \\
			\hline
		\end{tabular}
	\end{center} 
	Xác định tính đúng sai của các mệnh đề sau
	\choiceTF  
	{Mốt của mẫu số liệu lớn hơn $1$}  
	{\True Nhóm chứa mốt của mẫu số liệu là $[0{,}75; 1{,}25)$}  
	{Trung vị của mẫu số liệu là $M_e = 1{,}039$}  
	{Số trung bình của mẫu số liệu trên là $1{,}14$}  
	\loigiai{  
		\begin{itemchoice}  
			\itemch \textbf{Sai}. Mốt của một tập dữ liệu là giá trị có tần số lớn nhất. Trong bảng số liệu, tần số lớn nhất là $32$, thuộc nhóm $[0{,}75; 1{,}25)$.
			\[M_0 = u_m + \frac{n_m - n_{m-1}}{(n_m - n_{m-1}) + (n_m - n_{m+1})} \cdot (u_{m+1} - u_m)
			\]
			\[=0{,}75+\dfrac{32-25}{32-25+32-14}\cdot(1{,}25-0{,}75)=0{,}89.\]
			\itemch \textbf{Đúng}. Vì nhóm có tần số lớn nhất là $[0{,}75; 1{,}25)$, nên đây chính là nhóm chứa mốt.  
			\itemch \textbf{Sai}. Trung vị $M_e$ được tính theo công thức 
			\[M_e = u_m + \frac{\dfrac{n}{2} - C}{n_m} \cdot (u_{m+1} - u_m). 
			\]
			Với:  
			$n = 87$, tổng số lần đo. \\
			Nhóm chứa trung vị là $[0,75; 1,25)$.  
			Áp dụng công thức
			\[
			M_e = 0{,}75 + \dfrac{\dfrac{87}{2} - 25}{32} \cdot (1{,}25-0{,}75) = \dfrac{133}{128}\approx 1{,}03906.
			\]
			\itemch \textbf{Sai}. Số trung bình được tính theo công thức
			\[
			\overline{x} = \frac{n_1c_1 + \ldots + n_kc_k}{n} 
			\]
			\begin{center}
				\begin{tabular}{|c|c|c|c|c|c|}
					\hline \textbf{Thời gian đo} &$[0{,}25; 0{,}75)$&$[0{,}75; 1{,}25)$&$[1{,}25;1{,}75)$&$[1{,}75; 2{,}25)$&$[2{,}25; 2{,}75)$\\
					\hline \textbf{Giá trị đại diện} &$0{,}5$ & $1$&$1{,}5$ &$2$ &$2{,}5$ \\
					\hline \textbf{Số lần} &$25$ & $32$&$14$ &$12$ &$4$ \\
					\hline
				\end{tabular}
			\end{center} 
			Số trung bình của mẫu số liệu
			\[
			\overline{x} =\dfrac{(25 \cdot 0{,}5) + (32 \cdot 1) + (14 \cdot 1{,}5) + (12 \cdot 2) + (4 \cdot 2{,}5)}{87} =\dfrac{99{,}5}{87} \approx 1{,}14.
			\]
		\end{itemchoice}  
	}  
\end{ex}

\begin{bt}%[1D5H2-3]
% nguon: Archive/SP_SauLuoi_2024-2025/Sp_GHKII/data/GHKII/1-TN-DS-TLN-TL-THPT-Edison-HaiPhong-GKII-NH24-25.tex #258
%[Dự án đề kiểm tra Toán 10 GHKII NH24-25- Quan Ón]%[THPT Edison - Hải Phòng]
	Trong một cuộc đua Marathon được tổ chức ở thành phố A người ta thống kê lại được như sau
	\begin{center}
		\begin{tabular}{|l|c|c|c|c|c|}
			\hline
			Thời gian & $[120;140)$ & $[140;160)$ & $[160;180)$ & $[180;200)$ & $[200;220)$ \\
			\hline
			Số người & $4$ & $6$ & $10$ & $15$ & $25$\\
			\hline
		\end{tabular}
	\end{center}
	Tính $Q_3 - Q_1$ của mẫu số liệu ghép nhóm trên.
	\shortans{38}
	\loigiai{
		Số người tham gia cuộc đua Marathon là
		$$ n = 4 + 6 + 10 + 15 + 25 = 60 \textrm{ (người).} $$
		Gọi $x_1$; $x_2$; \dots $x_{60}$ lần lượt là thời gian chạy của $60$ người tham gia cuộc đua Marathon được xếp theo thứ tự không giảm.\\
		Do $x_1,\dots,x_4 \in [120;140)$; $x_5, \dots, x_{10} \in [140;160)$; $x_{11},\dots, x_{20} \in [160;180)$; $x_{21}, \dots, x_{35} \in [180;200)$; $x_{36}, \dots, x_{60} \in [200;220)$ nên trung vị của mẫu số liệu là $\dfrac{1}{2}\left(x_{30} + x_{31}\right)$.\\
		Tứ phân vị thứ nhất của dãy số liệu $x_1$; $x_2$; \dots, $x_{60}$ là $\dfrac{1}{2}\left(x_{15} + x_{16}\right) \in [160;180)$. Do đó, tứ phân vị thứ nhất của mẫu số liệu ghép nhóm là
		$$ Q_1 = 160 + \dfrac{\dfrac{1\cdot 60}{4} - (4 + 6)}{10}\cdot \left(180 - 160\right) = 170. $$
		Tứ phân vị thứ hai của dãy số liệu $x_1$; $x_2$; \dots, $x_{60}$ là $\dfrac{1}{2}\left(x_{45} + x_{46}\right) \in [200;220)$. Do đó, tứ phân vị thứ hai của mẫu số liệu ghép nhóm là
		$$ Q_3 = 200 + \dfrac{\dfrac{3\cdot 60}{4} - (4 + 6 + 10 + 15)}{25}\cdot \left(220 - 200\right) = 208. $$
		Do đó $Q_3 - Q_1 = 208 - 170 = 38$.
	}
\end{bt}

\begin{bt}%[1D5H1-3]
% nguon: Archive/2024/San pham - To Latex Strong - Tex hoa Lop 12 moi - Dot 2 (Chuong 1 + 2 + 3)/San pham - To Latex Strong - Tex hoa Lop 12 moi - Dot 2 (Chuong 1 + 2 + 3)/Data/Chuong 3/Chan troi sang tao/To 5 - Dot 17 - Chuong 3 - Bai 2 - CTST.tex #1415
%[Tổ 5 - Đợt 17 - Chương 3 - Bài 2 - CTST]%[TRẦN CHÁNH NGHĨA]
	Thầy Toán thống kê lại điểm trung bình cuối năm của các học sinh lớp 11A ở bảng sau:\\
	\begin{tabular}{|c|c|c|c|c| c|}
		\hline
		Điểm trung bình & $[5;6)$  & $[6;7)$  & $[7;8)$ & $[8;9)$  & $[9;10)$  \\
		\hline
		Số học sinh & $3$  & $7$ & $9$ & $15$ & $6$   \\
		\hline
	\end{tabular}\\
	Tính số trung bình của mẫu số liệu ghép nhóm ?
	\shortans{$7{,}85$}
	\loigiai
	{
		Ta có bảng thống kê điểm trung bình theo giá trị đại diện:\\
		\begin{tabular}{|c|c|c|c|c| c|}
			\hline
			Điểm trung bình & $[5;6)$  & $[6;7)$  & $[7;8)$ & $[8;9)$  & $[9;10)$  \\
			\hline
			Giá trị đại diện & $5{,}5$  & $6{,}5$  & $7{,}5$  & $8{,}5$  & $9{,}5$     \\
			\hline
			Số học sinh & $3$  & $7$ & $9$ & $15$ & $6$   \\
			\hline
		\end{tabular}\\
		Cỡ mẫu là $n = 3 + 7 + 9 + 15 + 6 = 40$. \\
		Số trung bình của mẫu số liệu ghép nhóm là\\
		$\overline x  = \dfrac{{3\cdot 5{,}5 + 7\cdot 6{,}5 + 9\cdot 7{,}5 + 15\cdot 8{,}5 + 6\cdot 9{,}5}}{{40}} = 7{,}85$. 
		
	}
\end{bt}

\begin{ex}%[1D5H2-2]
% nguon: Archive/2024/Update-SP-BG10-11-3/SP-3/BG11-3/data_1C5/DeOT-ChuongV-DeOnTap-So2-BuiManhTien.tex #1424
%[24-25- Bài Giảng K10-K11, Bùi Mạnh Tiến]
	Cho mẫu số liệu ghép nhóm về thống kê thời gian hoàn thành (phút) một bài kiểm tra trực tuyến của $100$ học sinh, ta có bảng số liệu sau
	\begin{center}
		\begin{tabular}{|c|c|c|c|c|c|c|}
			\hline
			Thời gian (phút)&$[33;35)$&$[35;37)$&$[37;39)$&$[39;41)$&$[41;43)$&$[43;45)$\\
			\hline
			Số học sinh&$4$&$13$&$38$&$27$&$14$&$4$\\
			\hline
		\end{tabular}
	\end{center}
	\choiceTF[t]
	{Thời gian hoàn thành bài trung bình của $100$ học sinh lớn hơn $39$ phút}
	{Trung vị của mẫu số liệu thuộc nhóm $[39;41)$}
	{\True Mốt của mẫu số liệu có giá trị lớn hơn $38$}
	{\True Trung vị của mẫu số liệu (làm tròn đến hàng phần trăm) là $38{,}74$}	
	\loigiai{
		\begin{itemchoice}
			\itemch Thời gian hoàn thành bài trung bình của $100$ học sinh là
			\begin{align*}
				\overline{x}=\dfrac{34\cdot 4+36\cdot 13+38\cdot 38+40\cdot 27+42\cdot 14+44\cdot 4}{100}=38{,}92 \text{ (phút)}.
			\end{align*}
			\itemch Vì $x_{50}$, $x_{51}$ thuộc nhóm $[37;39)$ nên $M_e$ thuộc nhóm này.
			\itemch Mốt của mẫu số liệu thuộc nhóm $[37;39)$. Ta có
			\begin{align*}
				M_0=37+\dfrac{38-13}{(38-13)+(38-27)}\cdot (39-37)\approx 38{,}4>38.
			\end{align*}
			\itemch Trung vị của mẫu số liệu (làm tròn đến hàng phần trăm) là
			\begin{align*}
				M_e=37+\dfrac{\dfrac{100}{2}-(4+13)}{38}\cdot (39-37)\approx 38{,}74.
			\end{align*}
		\end{itemchoice}
	}
\end{ex}

\begin{ex}%[1D5H2-3]
% nguon: Archive/2024/Update-SP-BG10-11-3/SP-3/BG11-3/data_1C5/DeOT-ChuongV-DeOnTap-So2-BuiManhTien.tex #1425
%[24-25- Bài Giảng K10-K11, Bùi Mạnh Tiến]
	Một cửa hàng trang sức khảo sát một số khách hàng xem họ dự định mua trang sức với mức giá nào (đơn vị: triệu đồng). Kết quả khảo sát được ghi lại ở bảng sau
	\begin{center}
		\begin{tabular}{|c|c|c|c|c|c|}
			\hline
			Mức giá&$[6;9)$&$[9;12)$&$[12;15)$&$[15;18)$&$[18;21)$\\
			\hline
			Số khách hàng&$20$&$75$&$48$&$23$&$12$\\
			\hline
		\end{tabular}
	\end{center}
	\choiceTF[t]
	{Cỡ mẫu của mẫu số liệu trên là $180$}
	{\True Giá trị trung bình của mẫu số liệu trên là $12{,}35$ (làm tròn đến hàng phần trăm)}
	{\True $Q_1=\dfrac{499}{50}$}
	{$Q_3=\dfrac{115}{8}$}	
	\loigiai{
		\begin{itemchoice}
			\itemch Cỡ mẫu của mẫu số liệu ghép nhóm trên là $n=20+75+48+23+12=178$.
			\itemch Giá trị trung bình của mẫu số liệu trên là
			\begin{align*}
				\overline{x}=\dfrac{7{,}5\cdot 20+10{,}5\cdot 75+13{,}5\cdot 48+16{,}5\cdot 23+19{,}5\cdot 12}{178}\approx 12{,}35.
			\end{align*}
			\itemch Nhóm chứa tứ phân vị thứ nhất là nhóm $[9;12)$, suy ra 
			\begin{align*}
				Q_1=a_2+\dfrac{\dfrac{n}{4}-m_1}{m_2}(a_3-a_2)=9+\dfrac{44{,}5-20}{75}\cdot \left(12-9\right)=\dfrac{499}{50}.
			\end{align*}
			\itemch Nhóm chứa tứ phân vị thứ ba là nhóm $[12;15)$, ta có $\dfrac{3n}{4}=133{,}5$, do đó
			\begin{align*}
				Q_3=a_3+\dfrac{\dfrac{3n}{4}-(m_1+m_2)}{m_3}\cdot \left(a_4-a_3\right)=12+\dfrac{133{,}5-95}{48}\cdot (15-12)=\dfrac{461}{32}.
			\end{align*}		
		\end{itemchoice}
	}
\end{ex}

\begin{bt}%[1D5H2-2]
% nguon: Archive/2024/Update-SP-BG10-11-3/SP-3/BG11-3/data_1C5/DeOT-ChuongV-DeOnTap-So2-BuiManhTien.tex #1428
%[Mức ]%[24-25- Bài Giảng K10-K11, Bùi Mạnh Tiến]
	Thời gian truy cập Internet mỗi buổi tối của một số học sinh được cho trong bảng sau
	\begin{center}
		\begin{tabular}{|c|c|c|c|c|c|}
			\hline
			Thời gian (phút)&$[10;20)$&$[20;30)$&$[30;40)$&$[40;50)$&$[50;60)$\\
			\hline
			Số học sinh&$3$&$12$&$15$&$24$&$2$\\
			\hline
		\end{tabular}
	\end{center}
	Nhóm chứa trung vị của mẫu số liệu trên có tần số là
	\shortans{$15$}
	\loigiai{
		Cỡ mẫu là $n=3+12+15+24+2=56$.\\
		Gọi $x_1$, $x_2$, $\ldots$; $x_{56}$ là thời gian truy cập Internet của các bạn học sinh được xếp theo thứ tự tăng dần.\\
		Trung vị là $\dfrac{x_{28}+x_{29}}{2}$. Nhóm chứa trung vị là nhóm $[30;40)$ có tần số là $15$.
	}
\end{bt}

\begin{bt}%[1D5H1-4]
% nguon: Archive/2024/Update-SP-BG10-11-3/SP-3/BG11-3/data_1C5/DeOT-ChuongV-DeOnTap-So2-BuiManhTien.tex #1429
%[Mức ]%[24-25- Bài Giảng K10-K11, Bùi Mạnh Tiến]
	Khảo sát thời gian chơi thể thao trong một ngày của $40$ học sinh lớp $11$B$3$ giáo viên thu được một mẫu dữ liệu ghép nhóm như sau
	\begin{center}
		\begin{tabular}{|c|c|c|c|c|c|c|}
			\hline
			Thời gian (phút)&$[30;40)$&$[40;50)$&$[50;60)$&$[60;70)$&$[70;80)$&$[80;90)$\\
			\hline
			Số học sinh&$2$&$10$&$16$&$8$&$2$&$2$\\
			\hline
		\end{tabular}
	\end{center}
	Mốt của mẫu số liệu trên là bao nhiêu? (làm tròn đến hàng phần mười).
	\shortans{$54{,}3$}
	\loigiai{
		Nhóm chứa mốt của mẫu số liệu trên là nhóm $[50;60)$.\\
		Do đó $u_m=50$, $n_{m-1}=10$, $n_m=16$; $n_{m+1}=8$; $u_{m+1}-u_m=60-50=10$.\\
		Mốt của mẫu số liệu ghép nhóm là
		\begin{align*}
			M_0=50+\dfrac{16-10}{(16-10)+(16-8)}\cdot 10\approx 54{,}3. 
		\end{align*}
	}
\end{bt}

\begin{bt}%[1D5H2-3]
% nguon: Archive/2024/Update-SP-BG10-11-3/SP-3/BG11-3/data_1C5/C5-B1-Dang4-5-PhanAnh.tex #1434
[Mức độ 2]%[Bài Giảng 10-11 - Phan Anh]
	\immini
	{
		Bảng bên cho biết tần số ghép nhóm số liệu thống kê cân nặng của $40$ học sinh lớp $11A$ trong một trường trung học phổ thông (đơn vị: kilôgam). Xác định tứ phân vị của mẫu số liệu ghép nhóm.
	}
	{
		\begin{tabular}{|c|c|c|}
			\hline
			\textbf{Nhóm} &\textbf{Tần số} &\textbf{Tần số tích luỹ}\\ 
			\hline
			$\left[30;40\right)$ & $2$ & $2$\\
			$\left[40;50\right)$ & $10$ & $12$\\
			$\left[50;60\right)$ & $16$ & $28$\\
			$\left[60;70\right)$ & $8$ & $36$\\
			$\left[70;80\right)$ & $2$ & $38$\\
			$\left[80;90\right)$ & $2$ & $40$\\
			\hline
			& $n=40$ &\\
			\hline
		\end{tabular}
	}
	\loigiai{
		Số phần tử của mẫu là $n=40$.
		\begin{itemize}
			\item Ta có $\dfrac{n}{4}=10$, mà $2<10<12$ nên nhóm $2$ là nhóm đầu tiên có tần số tích luỹ lớn hơn hoặc bằng $10$.\\
			Xét nhóm $2$ là nhóm $\left[40;50\right)$ có $s=40$, $h=10$, $n_2=10$ và nhóm $1$ có $cf_1=2$.\\
			Khi đó, tứ phân vị thứ nhất là
			\[
			Q_1=40 +\dfrac{10-2}{10}\cdot 10=48\text{ (kg)}.
			\]
			\item Ta có $\dfrac{n}{2}=20$, mà $12<20<28$ nên nhóm $3$ là nhóm đầu tiên có tần số tích luỹ lớn hơn hoặc bằng $20$.\\
			Xét nhóm $3$ là nhóm $\left[50;60\right)$ có $r=50$, $d=10$, $n_3=16$ và nhóm $2$ có $cf_2=12$.\\
			Khi đó, tứ phân vị thứ hai là
			\[
			Q_2=50 +\dfrac{20-12}{16}\cdot 10=55\text{ (kg)}.
			\]
			\item Ta có $\dfrac{3n}{4}=30$, mà $28<30<36$ nên nhóm $4$ là nhóm đầu tiên có tần số tích luỹ lớn hơn hoặc bằng $30$.\\
			Xét nhóm $4$ là nhóm $\left[60;70\right)$ có $t=50$, $l=10$, $n_4=8$ và nhóm $3$ có $cf_3=28$.\\
			Khi đó, tứ phân vị thứ ba là
			\[
			Q_3=60 +\dfrac{30-28}{8}\cdot 10=62{,}5\text{ (kg)}.
			\]
		\end{itemize}
		Vậy tứ phân vị của mẫu số liệu trên là $48$, $55$ và $62{,}5$.
	}
\end{bt}

\begin{ex}%[1D5H2-3]
% nguon: Archive/2024/Update-SP-BG10-11-3/SP-3/BG11-3/data_1C5/C5-B1-Dang4-5-PhanAnh.tex #1435
[Mức độ 2]%[Bài Giảng 10-11 - Phan Anh]
	Khảo sát thời gian xem điện thoại trong một ngày của một số học sinh khối $12$ thu được mẫu số liệu ghép nhóm sau
	\begin{center}
		\renewcommand{\arraystretch}{1.2}
		\begin{tabular}{|c|c|c|c|c|c|}
			\hline Thời gian (phút) & $[0;20)$ & $[20;40)$ & $[40;60)$ & $[60;80)$ & $[80;100)$\\
			\hline Số học sinh & $4$ & $8$ & $12$ & $10$ & $8$\\
			\hline
		\end{tabular}
	\end{center}
	Các khẳng định sau đây đúng hay sai?
	\choiceTF[t]
	{\True Tổng số học sinh được khảo sát là $42$}
	{Số trung bình của mẫu số liệu lớn hơn $54{,}8$}
	{\True Tứ phân vị thứ nhất lớn hơn $35$}
	{\True Tứ phân vị thứ ba nhỏ hơn $76$}
	\loigiai{
		\begin{itemchoice}
			\itemch Tổng số học sinh được khảo sát là $n=4+8+12+10+8=42$.
			\itemch Số trung bình của mẫu số liệu là 
			\[\overline{x}=\dfrac{4\cdot10+8\cdot30+12\cdot50+10\cdot70+8\cdot90}{4+8+12+10+8}=\dfrac{1150}{21}\approx 54{,}76.\]
			\itemch Gọi $x_1$, $x_2$, \ldots, ${x_{42}}$ là thời gian xem điện thoại trong ngày của $42$ học sinh khối $12$ và giả sử dãy này đã sắp xếp theo thứ tự tăng dần.\\
			Khi đó tứ phân vị thứ nhất $Q_1$ là trung vị của dãy $x_1$, $x_2$,..., $x_{21}$ nên $Q_1=x_{11}$. Do đó $Q_1$ thuộc nhóm $[20;40)$. \\
			Suy ra $Q_1=20 +\dfrac{\dfrac{42}{4} - 4}{8}\cdot (40-20)=36{,}25$.
			\itemch 
			Tứ phân vị thứ ba $Q_3$ là trung vị của dãy $x_{22}$, $x_{23}$,..., $x_{42}$ nên $Q_3=x_{32}$. Do đó $Q_3$ thuộc nhóm $[60;80)$.\\
			Ta có $Q_3=60 +\dfrac{\dfrac{3\cdot 42}{4} - 24}{10}\cdot (80-60)=75$. 
		\end{itemchoice}
	}
\end{ex}

\begin{bt}%[1D5H2-3]
% nguon: Archive/2024/Update-SP-BG10-11-3/SP-3/BG11-3/data_1C5/C5-B1-Dang4-5-PhanAnh.tex #1436
[Mức độ 2]%[Bài Giảng 10-11 - Phan Anh]
	Thời gian (phút) truy cập internet mỗi buổi tối của một số học sinh được cho trong bảng sau
	\begin{center}
		\begin{tabular}{|c|c|c|c|c|c|c|}
			\hline
			Thời gian (phút)	& $\left[9{,}5;12{,}5\right)$ & $\left[12{,}5;15{,}5\right)$ & $\left[15{,}5;18{,}5\right)$ & $\left[18{,}5;21{,}5\right)$ & $\left[21{,}5;24{,}5\right)$\\
			\hline
			Số học sinh&$3$	& $12$ & $15$ &$24$ & $2$\\
			\hline
		\end{tabular}	
	\end{center}
	Tìm tứ phân vị thứ nhất $Q_1$ và tứ phân vị thứ ba $Q_3$ của mẫu số liệu ghép nhóm.
	\loigiai{
		Cỡ mẫu là $n=3+12+15+24+2=56$.\\
		Gọi $x_1$, $x_2$, \ldots, $x_{56}$ là thời gian (phút) truy cập internet mỗi buổi tối của $56$ học sinh được xếp theo thứ tự không giảm.\\
		Tứ phân vị thứ nhất $Q_1$ là $\dfrac{x_{14}+x_{15}}{2}$. Do $2$ giá trị $x_{14}$, $x_{15}$ thuộc nhóm $\left[12{,}5;15{,}5\right)$ nên nhóm này chứa $Q_1$. Do đó, $p=2$;$a_2=12{,}5$;$m_2=12$;$m_1=3$;$a_3-a_2=3$ và ta có $$Q_1=12{,}5+\dfrac{\dfrac{56}{4}-3}{12}\cdot 3=15{,}25.$$
		Với tứ phân vị thứ ba $Q_3$ là $\dfrac{x_{42}+x_{43}}{2}$. Do $2$ giá trị $x_{42},\,x_{43}$ thuộc nhóm $\left[18{,}5;21{,}5\right)$ nên nhóm này chứa $Q_3$. Do đó, $p=4$;$a_4=18{,}5$;$m_4=24$;$m_1+m_2+m_3=3+12+15=30$;$a_5-a_4=3$ và ta có $$Q_3=18{,}5+\dfrac{\dfrac{3\cdot 56}{4}-30}{24}\cdot 3=20.$$
	}
\end{bt}

\begin{bt}%[1D5H2-3]
% nguon: Archive/2024/Update-SP-BG10-11-3/SP-3/BG11-3/data_1C5/C5-B1-Dang4-5-PhanAnh.tex #1437
[Mức độ 2]%[Bài Giảng 10-11 - Phan Anh]
	Điểm thi môn Toán (thang điểm 100, điểm được làm tròn đến 1) của $60$ thí sinh được cho trong bảng sau
	\begin{center}
		\begin{tabular}{|c|c|c|c|c|c|}
			\hline
			Điểm	& $0-9$ & $10-19$ & $20-29$ & $30-39$ & $40-49$\\
			\hline
			Số thí sinh	& $1$ & $2$ & $4$ & $6$ & $15$\\
			\hline
			Điểm	& $50-59$ & $60-69$ & $70-79$ & $80-89$ & $90-99$\\
			\hline
			Số thí sinh	& $12$ & $10$ & $6$ & $3$ & $1$\\
			\hline
		\end{tabular}
	\end{center}
	\begin{enumerate}
		\item Hiệu chỉnh để thu được mẫu số liệu ghép nhóm với độ dài mỗi nhóm là $20$.
		\item Tìm các tứ phân vị và giải thích ý nghĩa của chúng.
	\end{enumerate}
	\loigiai{
		\begin{enumerate}
			\item Bảng số liệu ghép nhóm về điểm thi môn Toán của 60 thí sinh
			\begin{center}
				\begin{tabular}{|c|c|c|c|c|c|}
					\hline
					Điểm	& $\left[0;20\right)$ & $\left[20;40\right)$ & $\left[40;60\right)$ & $\left[60;80\right)$ & $\left[80;100\right)$\\
					\hline
					Số thí sinh	& $3$ & $10$ & $27$ & $16$ & $4$\\
					\hline
				\end{tabular}
			\end{center}
			\item Cỡ mẫu $n=60$. Gọi $x_1$, $x_2$,$\ldots$, $x_{60}$ là điểm thi môn Toán của 60 học sinh và giả sử dãy này đã được sắp xếp theo thứ tự tăng dần. Khi đó, trung vị là $\dfrac{x_{30}+x_{31}}{2}$.\\
			Do hai giá trị $x_{30}$, $x_{31}$ thuộc nhóm $\left[40;60\right)$ nên nhóm này chứa trung vị. Do đó, $p=3;a_3=40;m_3=27;m_1+m_2=13;a_4-a_3=20$ và ta có
			$$Q_2=M_e=40+\dfrac{\dfrac{60}{2}-13}{27}\cdot 20\approx 52{,}6.$$
			Tứ phân vị thứ nhất $Q_1=\dfrac{x_{15}+x_{16}}{2}$. Do hai giá trị $x_{15}$, $x_{16}$ thuộc nhóm $\left[40;60\right)$ nên nhóm này chứa $Q_1$. Do đó, $p=3;\,a_3=40;\,m_3=27;\,m_1+m_2=13;\,a_4-a_3=20$ và ta có
			$$Q_1=40+\dfrac{\dfrac{60}{4}-13}{27}\cdot 20\approx 41{,}5.$$
			Tứ phân vị thứ ba $Q_3=\dfrac{x_{45}+x_{46}}{2}$. Do hai giá trị $x_{45}$, $x_{46}$ thuộc nhóm $\left[60;80\right)$ nên nhóm này chứa $Q_3$. Do đó, $p=4;\,a_4=60;\,m_4=16;\,m_1+m_2+m_3=40;\,a_5-a_4=20$ và ta có
			$$Q_3=60+\dfrac{\dfrac{3\cdot 60}{4}-40}{16}\cdot 20\approx 66{,}3.$$
			Khoảng cách từ $Q_1$ đến $Q_2$ là $11{,}1$ còn khoảng cách từ $Q_2$ và $Q_3$ là $13{,}7$. Điều này cho thấy mẫu số liệu tập trung với mật độ cao hơn ở bên trái $Q_2$ và mật độ thấp hơn ở bên phải $Q_2$.
		\end{enumerate}	
	}
\end{bt}

\begin{bt}%[1D5V2-3]
% nguon: Archive/2024/Update-SP-BG10-11-3/SP-3/BG11-3/data_1C5/C5-B1-Dang4-5-PhanAnh.tex #1438
[Mức độ 3]%[Bài Giảng 10-11 - Phan Anh]
	\immini{
		Phỏng vấn một số học sinh lớp $11$ về thời gian (giờ) ngủ của một buổi tối, thu được bảng số liệu ở bên.
		\begin{enumerate}
			\item So sánh thời gian ngủ trung bình của các bạn học sinh nam và nữ.
			\item Hãy cho biết $75\%$ học sinh khối $11$ ngủ ít nhất bao nhiêu giờ?
		\end{enumerate}
	}
	{
		\begin{tabular}{|c|c|c|}
			\hline
			Thời gian	& Số học sinh nam & Số học sinh nữ\\
			\hline
			$\left[4;5\right)$	& $6$ & $4$\\
			\hline
			$\left[5;6\right)$	& $10$ & $8$\\
			\hline
			$\left[6;7\right)$	& $13$ & $10$\\
			\hline
			$\left[7;8\right)$	& $9$ & $11$\\
			\hline
			$\left[8;9\right)$	& $7$ & $8$\\
			\hline
		\end{tabular}
	}
	\loigiai{
		\begin{enumerate}
			\item Trong mỗi khoảng thời gian, giá trị đại diện là trung bình cộng của giá trị hai đầu mút nên ta có bảng sau
			\begin{center}
				\begin{tabular}{|c|c|c|}
					\hline
					Thời gian	& Số học sinh nam & Số học sinh nữ\\
					\hline
					$4{,}5$	& $6$ & $4$\\
					\hline
					$5{,}5$	& $10$ & $8$\\
					\hline
					$6{,}5$	& $13$ & $10$\\
					\hline
					$7{,}5$	& $9$ & $11$\\
					\hline
					$8{,}5$	& $7$ & $8$\\
					\hline
				\end{tabular}	
			\end{center}
			Tổng số học sinh nam là $n_1=6+10+13+9+7=45$.\\ Thời gian ngủ trung bình của học sinh nam là
			$$\overline{x}_1=\dfrac{4{,}5\cdot 6+5{,}5\cdot10+6{,}5\cdot13+7{,}5\cdot9+8{,}5\cdot7}{45}=\dfrac{587}{90}\approx 6{,}52\,\,\text{(giờ)}.$$
			Tổng số học sinh nữ là $n_2=4+8+10+11+8=41$. Thời gian ngủ trung bình của học sinh nữ là
			$$\overline{x}_2=\dfrac{4,5\cdot4+5,5\cdot8+6,5\cdot10+7,5\cdot11+8,5\cdot8}{41}=\dfrac{555}{82}\approx 6{,}77\,\,\text{(giờ)}.$$
			Vì $\overline{x}_2>\overline{x}_1$ nên thời gian ngủ trung bình của các bạn học sinh nữ lớn hơn thời gian ngủ trung bình của các bạn nam.
			\item Tổng số học sinh được điều tra là $n=n_1+n_2=45+41=86$.\\
			Giả sử $x_1$; $x_2$; $x_3$; \ldots; $x_{86}$ là dãy giá trị được sắp xếp theo thứ tự không giảm.\\
			Ta có bảng sau
			\begin{center}
				\begin{tabular}{|c|c|c|}
					\hline
					Thời gian	& Số học sinh\\
					\hline
					$\left[4;5\right)$	& $10$\\
					\hline
					$\left[5;6\right)$	& $18$\\
					\hline
					$\left[6;7\right)$	& $23$\\
					\hline
					$\left[7;8\right)$	& $20$\\
					\hline
					$\left[8;9\right)$	& $15$\\
					\hline
				\end{tabular}
			\end{center}
			Tứ phân vị thứ nhất $Q_1$ là $x_{22}$. Do $x_{22}$ thuộc nhóm $\left[5;6\right)$ nên nhóm này chứa $Q_1$.\\ Do đó, $p=2;\,a_2=5;\,m_2=18;\,m_1=10;\,a_3-a_2=1$ và ta có
			$$Q_1=5+\dfrac{\dfrac{86}{4}-10}{18}\cdot 1=\dfrac{203}{36}\approx 5{,}64\,\text{(giờ)}.$$
			Nghĩa là có $25\%$ học sinh khối $11$ ngủ ít hơn $5{,}64$ giờ.\\ Vậy $75\%$ học sinh khối $11$ ngủ ít nhất $5{,}64$ giờ.
		\end{enumerate}
	}
\end{bt}

\begin{bt}%[1D5H2-3]
% nguon: Archive/2024/Update-SP-BG10-11-3/SP-3/BG11-3/data_1C5/C5-B1-Dang4-5-PhanAnh.tex #1439
[Mức độ 2]%[Bài Giảng 10-11 - Phan Anh]
	Số điểm một cầu thủ bóng rổ ghi được trong $ 20 $ trận đấu được cho ở bảng sau
	\begin{center}
		\begin{tabular}{|c|c|c|c|c|c|c|c|c|c|c|c|c|c|c|c|c|c|c|c|}
			\hline 
			$ 25 $	& $ 23 $ & $ 21 $ & $ 13 $ & $ 8 $ & $ 14 $ & $ 15 $ & $ 18 $ & $ 22 $ & $ 11 $\\ 
			\hline 
			$ 24 $	& $ 12 $ & $ 14 $ & $ 14 $ & $ 18 $ & $ 6 $ & $ 8 $ & $ 25 $ & $ 10 $ & $ 11 $ \\ 
			\hline 
		\end{tabular} 
	\end{center}
	\begin{enumerate}
		\item Tìm tứ phân vị của dãy số liệu trên.
		\item Tổng hợp lại dãy số liệu trên vào bảng tần số ghép nhóm theo mẫu sau
		\begin{center}
			\begin{tabular}{|c|c|c|c|c|}
				\hline 
				\textbf{Điểm số}	& $\left[ 6;10\right] $ & $\left[ 11;15\right] $ & $\left[16;20\right] $ & $\left[ 21;25\right] $\\ 
				\hline 
				\textbf{Số trận}	& ? & ? & ? & ?\\ 
				\hline 
			\end{tabular} 
		\end{center}
		\item Hãy ước lượng tứ phân vị của số liệu từ bảng tần số ghép nhóm trên.
	\end{enumerate}
	\loigiai{
		\begin{enumerate}
			\item Gọi $ x_1 $, $ x_2 $,\ldots, $ x_{20} $ là mẫu số liệu được xếp theo thứ tự không giảm. Ta có bảng sau
			\begin{center}
				\begin{tabular}{|c|c|c|c|c|c|c|c|c|c|}
					\hline 
					$ x_1 $	& $ x_2 $ & $ x_3 $ & $ x_4 $ & $ x_5 $ & $ x_6 $ & $ x_7 $ & $ x_8 $ & $ x_9 $ & $ x_{10} $\\ 
					\hline 
					$ 6 $	& $ 8 $ & $ 8 $ & $ 10 $ & $ 11 $ & $ 11 $ & $ 12 $ & $ 13 $ & $ 14 $ & $ 14 $\\ 
					\hline 
					$ x_{11} $	& $ x_{12} $ & $ x_{13} $ & $ x_{14} $ & $ x_{15} $ & $ x_{16} $ & $ x_{17} $ & $ x_{18} $ & $ x_{19} $ & $ x_{20} $\\ 
					\hline 
					$ 14 $	& $ 15 $ & $ 18 $ & $ 18 $ & $ 21 $ & $ 22 $ & $ 23 $ & $ 24 $ & $ 25 $ & $ 25 $\\
					\hline 
				\end{tabular} 
			\end{center}
			Tứ phân vị thứ nhất của mẫu số liệu là $\dfrac{1}{2}\cdot\left( x_{5}+x_6\right) =\dfrac{1}{2}\cdot\left( 11+11\right)=11$.\\
			Tứ phân vị thứ hai của mẫu số liệu là $\dfrac{1}{2}\cdot\left( x_{10}+x_{11}\right)=\dfrac{1}{2}\cdot\left( 14+14\right)=14$.\\
			Tứ phân vị thứ ba của mẫu số liệu là $\dfrac{1}{2}\cdot\left( x_{15}+x_{16}\right)=\dfrac{1}{2}\cdot\left( 21+22\right)=21{,}5$.
			\item Tổng hợp lại dãy số liệu trên vào bảng tần số ghép nhóm đã cho, ta được
			\begin{center}
				\begin{tabular}{|c|c|c|c|c|}
					\hline 
					\textbf{Điểm số}	& $\left[ 6;10\right] $ & $\left[ 11;15\right] $ & $\left[16;20\right] $ & $\left[ 21;25\right] $\\ 
					\hline 
					\textbf{Số trận}	& $ 4 $ & $ 8 $ & $ 2 $ & $ 6 $\\ 
					\hline 
				\end{tabular} 
			\end{center}
			\item Vì số điểm của cầu thủ bóng rổ là số nguyên nên ta hiệu chỉnh bảng số liệu tần số ghép nhóm như sau
			\begin{center}
				\begin{tabular}{|c|c|c|c|c|}
					\hline 
					\textbf{Điểm số}	& $\left[ 5{,}5;10{,}5\right) $ & $\left[ 10{,}5;15{,}5\right) $ & $\left[15{,}5;20{,}5\right) $ & $\left[ 20{,}5;25{,}5\right) $\\ 
					\hline 
					\textbf{Số trận}	& $ 4 $ & $ 8 $ & $ 2 $ & $ 6 $\\ 
					\hline 
				\end{tabular} 
			\end{center}
			Gọi $ x_1 $, $ x_2 $, $\ldots $, $ x_{20} $ là mẫu số liệu được xếp theo thứ tự không giảm.\\
			Ta có $ x_1,\ldots, x_4\in \left[ 5{,}5;10{,}5\right) $;$ x_5$, \ldots, $x_{12}\in \left[ 10{,}5;15{,}5\right)$;$ x_{13}, x_{14}\in \left[15{,}5;20{,}5\right) $;$ x_{15},\ldots, x_{20}\in \left[ 20{,}5;25{,}5\right) $.\\
			Tứ phân vị thứ nhất của mẫu số liệu là $\dfrac{1}{2}\cdot\left( x_{5}+x_6\right)\in\left[ 10{,}5;15{,}5\right)$ nên tứ phân vị thứ nhất của mẫu số liệu ghép nhóm là $$Q_1=10{,}5 +\dfrac{\dfrac{20}{4}-4}{8}\cdot\left(15{,}5-10{,}5\right)=\dfrac{89}{8}=11{,}125.$$
			Tứ phân vị thứ hai của mẫu số liệu là $\dfrac{1}{2}\cdot\left( x_{10}+x_{11}\right)\in\left[ 10{,}5;15{,}5\right)$ nên tứ phân vị thứ hai của mẫu số liệu ghép nhóm là $$Q_2=10{,}5 +\dfrac{\dfrac{20}{2}-4}{8}\cdot\left(15{,}5-10{,}5\right)=\dfrac{57}{4}=14{,}25.$$
			Tứ phân vị thứ ba của mẫu số liệu là $\dfrac{1}{2}\cdot\left( x_{15}+x_{16}\right)\in\left[ 20{,}5;25{,}5\right)$ nên tứ phân vị thứ ba của mẫu số liệu ghép nhóm là $$Q_3=20{,}5 +\dfrac{\dfrac{20\cdot 3}{4}-(4+8+2)}{6}\cdot\left(25{,}5-20{,}5\right)=\dfrac{64}{3}\approx 21{,}333.$$
		\end{enumerate}
	}
\end{bt}

\begin{ex}%[1D5H2-3]
% nguon: Archive/2024/Update-SP-BG10-11-3/SP-3/BG11-3/data_1C5/C5-B1-Dang4-5-PhanAnh.tex #1440
%[Bài Giảng 10-11 - Phan Anh]
	Thời gian (phút) xem tivi mỗi buổi tối của một số học sinh được cho trong bảng sau
	\begin{center}
		\begin{tabular}{|c|c|c|c|c|c|c|c|}
			\hline
			{Thời gian}&{$[6{,}5;9{,}5)$} & {$[9{,}5;12{,}5)$} &{$[12{,}5;15{,}5)$}&{$[15{,}5;18{,}5)$} &{$[18{,}5;21{,}5)$}&{$[21{,}5;24{,}5)$}&{$[24{,}5;27{,}5)$}\\
			\hline
			{Số học sinh}&{$2$} & {$3$} &{$12$}&{$15$} &{$24$}&{$2$}&{$2$}\\
			\hline
		\end{tabular}
	\end{center}
	\choiceTF[t]
	{Giá trị đại diện của nhóm $[12{,}5;15{,}5)$ là $14{,}5$}
	{Giá trị trung bình của mẫu số liệu là $17$}
	{\True Trung vị của mẫu số liệu là $18{,}1$}
	{\True Các tứ phân vị của mẫu số liệu là $Q_1=15$, $Q_2=18{,}1$, $Q_3=20{,}125$}
	\loigiai{
		\begin{itemchoice}
			\itemch
			Giá trị đại diện của nhóm $[12{,}5;15{,}5)$ là $\dfrac{12{,}5+15{,}5}{2}=14$.
			\itemch
			Bảng giá trị đại diện
			\begin{center}
				\begin{tabular}{|c|c|c|c|c|c|c|c|}
					\hline
					{Thời gian}&{$[6{,}5;9{,}5)$} & {$[9{,}5;12{,}5)$} &{$[12{,}5;15{,}5)$}&{$[15{,}5;18{,}5)$} &{$[18{,}5;21{,}5)$}&{$[21{,}5;24{,}5)$}&{$[24{,}5;27{,}5)$}\\
					\hline
					{Số HS}&{$2$} & {$3$} &{$12$}&{$15$} &{$24$}&{$2$}&{$2$}\\
					\hline
					{GTĐD}& $8$ & $11$ & $14$ & $17$ & $20$ & $23$ & $26$\\
					\hline
				\end{tabular}
			\end{center}
			Cỡ mẫu $n=60$.\\
			Giá trị trung bình của mẫu số liệu
			\[\overline{x}=\dfrac{2\cdot 8+3\cdot 11+12\cdot 14+15\cdot 17+24\cdot 20+2\cdot 23+2\cdot 26}{60}=17{,}5.\]
			\itemch
			Giả sử $x_1$, $x_2$,\ldots, $x_{60}$ là các giá trị gốc của mẫu số liệu đã được xếp theo thứ tự không giảm.\\
			Do cỡ mẫu $n=60$ nên giá trị trung vị là $M_e=\dfrac{x_{30}+x_{31}}{2}\in [15{,}5;18{,}5)$.\\
			Số trung vị của mẫu số liệu ghép nhóm này là
			\[M_e=15{,}5+\dfrac{\dfrac{60}{2}-(2+3+12)}{15}\cdot (18{,}5-15{,}5)=18{,}1.\]
			\itemch
			Giả sử $x_1$, $x_2$,\ldots, $x_{60}$ là các giá trị gốc của mẫu số liệu đã được xếp theo thứ tự không giảm.\\
			Do cỡ mẫu $n=60$ nên
			\begin{eqnarray*}
				Q_1=\dfrac{x_{15}+x_{16}}{2}\in [12{,}5;15{,}5),\, Q_2=\dfrac{x_{30}+x_{31}}{2}\in [15{,}5;18{,}5)\text{ và } Q_3=\dfrac{x_{45}+x_{46}}{2}\in [18{,}5;21{,}5).
			\end{eqnarray*}
			Áp dụng công thức tính cho các tứ phân vị ta có
			\allowdisplaybreaks
			\begin{eqnarray*}
				&&Q_1=12{,}5 +\dfrac{\dfrac{60}{4}-(2+3)}{12}\cdot (15{,}5-12{,}5)=15;\\
				&&Q_2=15{,}5+\dfrac{\dfrac{60}{2}-(2+3+12)}{15}\cdot (18{,}5-15{,}5)=18{,}1;\\
				&&Q_3=18{,}5+\dfrac{3\cdot\dfrac{60}{4}-(2+3+12)}{15}\cdot (21{,}5-18{,}5)=20{,}125.
			\end{eqnarray*}
		\end{itemchoice}
	}
\end{ex}

\begin{ex}%[1D5V2-3]
% nguon: Archive/2024/Update-SP-BG10-11-3/SP-3/BG11-3/data_1C5/C5-B1-Dang4-5-PhanAnh.tex #1441
%[Bài Giảng 10-11 - Phan Anh]
	Phỏng vấn một số học sinh khối $11$ về thời gian (giờ) ngủ của một buổi tối, thu được bảng số liệu sau
	\begin{center}
		\renewcommand{\arraystretch}{1.5}
		\begin{tabular}{|c|c|c|}
			\hline
			Thời gian & Số học sinh nam & Số học sinh nữ\\
			\hline
			$[4;5)$ & $8$ & $4$\\
			\hline
			$[5;6)$ & $10$ & $6$\\
			\hline
			$[6;7)$ & $13$ & $12$\\
			\hline
			$[7;8)$ & $15$ & $18$\\
			\hline
			$[8;9)$ & $7$ & $8$\\
			\hline
		\end{tabular}
	\end{center}
	\choiceTF[t]
	{\True Độ dài của nhóm bằng $1$}
	{Thời gian ngủ trung bình của các bạn học sinh nam nhiều hơn các bạn học sinh nữ}
	{\True Phần lớn học sinh được khảo sát trong khối $11$ ngủ nhiều hơn $6{,}5$ giờ}
	{$75$\% học sinh được khảo sát trong khối $11$ ngủ ít nhất $5{,}5$ giờ}
	\loigiai{
		\begin{itemchoice}
			\itemch 
			Dựa vào bảng số liệu, ta thấy độ dài của mỗi nhóm là $1$.
			\itemch 
			Từ bảng thống kê ta có
			\begin{center}
				\renewcommand{\arraystretch}{1.5}
				\begin{tabular}{|c|c|c|c|}
					\hline
					Thời gian & Giá trị đại diện & Số học sinh nam & Số học sinh nữ\\
					\hline
					$[4;5)$ & $4{,}5$ & $8$ & $4$\\
					\hline
					$[5;6)$ & $5{,}5$ & $10$ & $6$\\
					\hline
					$[6;7)$ & $6{,}5$ & $13$ & $12$\\
					\hline
					$[7;8)$ & $7{,}5$ & $15$ & $18$\\
					\hline
					$[8;9)$ & $8{,}5$ & $7$ & $8$\\
					\hline
				\end{tabular}
			\end{center}
			Thời gian ngủ trung bình của các bạn nam là 
			$$\overline{x}_{\text {nam }}=\dfrac{4{,}5\cdot 8+5{,}5\cdot 10+6{,}5\cdot 13+7{,}5\cdot 15+8{,}5\cdot 7}{8+10+13+15+7}\approx 6{,}56.$$
			Thời gian ngủ trung bình của các bạn nữ là 
			$$\overline{x}_{\text {nữ} }=\dfrac{4{,}5\cdot 4+5{,}5\cdot 6+6{,}5\cdot 12+7{,}5\cdot 18+8{,}5\cdot 8}{4+6+12+18+8}\approx 6{,}92.$$
			Vì $6{,}56<6{,}92$ nên thời gian ngủ trung bình của các bạn nữ nhiều hơn các bạn nam.
			\itemch 
			Từ bảng thống kê ta có
			\begin{center}
				\renewcommand{\arraystretch}{1.5}
				\begin{tabular}{|c|c|c|c|c|}
					\hline
					Thời gian & Giá trị đại diện & Số học sinh nam & Số học sinh nữ & Số học sinh cả nam và nữ \\
					\hline
					$[4;5)$ & $4{,}5$ & $8$ & $4$ & $12$\\
					\hline
					$[5;6)$ & $5{,}5$ & $10$ & $6$ & $16$\\
					\hline
					$[6;7)$ & $6{,}5$ & $13$ & $12$ & $25$\\
					\hline
					$[7;8)$ & $7{,}5$ & $15$ & $18$ & $33$\\
					\hline
					$[8;9)$ & $8{,}5$ & $7$ & $8$ & $15$\\
					\hline
				\end{tabular}
			\end{center}
			Cỡ mẫu $n=101$.\\
			Gọi $x_1$, $x_2$, \ldots, $x_{101}$ là thời gian ngủ của $101$ học sinh khối 11 được xếp theo thứ tự không giảm.\\
			Trung vị của mẫu số liệu là $x_{51}$ thuộc nhóm $[6;7)$.\\
			Do đó, $p=3$, $ a_{3}=6$, $ a_{4}=7$, $ m_{2}=16$, $m_{1}=12$, $ a_{4}-a_{3}=1$ .\\
			Ta có trung vị của mẫu số liệu là
			$$
			M_{e}=6+\dfrac{\dfrac{101}{2}-(12+16)}{25}\cdot(7-6)=6{,}9.
			$$
			Vậy $50$\% học sinh có thời gian ngủ nhiều hơn $6{,}9$ giờ.
			\itemch 
			Cỡ mẫu $n=101$.\\
			Gọi $x_1$, $x_2$, \ldots, $x_{101}$ là thời gian ngủ của $101$ học sinh khối 11 được xếp theo thứ tự không giảm.\\
			Tứ phân vị thứ nhất $Q_{1}$ là $\dfrac{x_{25}+x_{26}}{2}$.\\
			Do $x_{25}$, $x_{26}$ đều thuộc nhóm $[5;6)$ nên nhóm này chứa $Q_{1}$.\\
			Do đó, $p=2$, $ a_{2}=5$, $ m_{2}=16$, $ m_{1}=12$, $ a_{3}-a_{2}=1$ và ta có\\
			$$Q_{1}=5+\dfrac{\dfrac{101}{4}-12}{16}\cdot 1\approx 5{,}83.$$
			Vậy $75\%$ học sinh khối $11$ ngủ ít nhất $5{,}83$ giờ.\\
			Mệnh đề sai.
		\end{itemchoice}
	}
\end{ex}

\begin{ex}%[1D5H2-3]
% nguon: Archive/2024/Update-SP-BG10-11-3/SP-3/BG11-3/data_1C5/C5-B1-Dang4-5-PhanAnh.tex #1442
%[Bài Giảng 10-11 - Phan Anh]
	Khảo sát cân nặng trung bình của học sinh lớp 11A cho trong bảng dưới đây
	\begin{center}
		\begin{tabular}{|c|c|c|c|c|c|c|}
			\hline
			Cân nặng (kg) &$[40{,}5;45{,}5)$ &$[45{,}5;50{,}5)$ &$[50{,}5;55{,}5)$ &$[55{,}5;60{,}5)$ &$[60{,}5;65{,}5)$
			&$[65{,}5;70{,}5)$ \\
			\hline
			Số học sinh	&$9$ &$7$ &$16$ &$5$ & $2$ & $3$ \\
			\hline
		\end{tabular}
	\end{center}
	Xét tính đúng - sai của các khẳng định sau
	\choiceTF[t]
	{\True Số học sinh được khảo sát là $42$}
	{\True Cân nặng trung bình của học sinh lớp 11A không ít hơn $52$ kg}
	{\True Nhóm chứa tứ trung vị là $[50{,}5;55{,}5)$}
	{Tứ phân vị thứ nhất lớn hơn $47$}
	\loigiai{
		\begin{itemchoice}
			\itemch 
			Cỡ mẫu là $n=9+7+16+5+2+3=42$.
			\itemch 
			Cân nặng trung bình của học sinh lớp 11A là
			$$\overline{x}=\dfrac{9\cdot 43+7\cdot 48+16\cdot 53+5\cdot 58+2\cdot 63+3\cdot 68}{42}\approx 52{,}17\,\text{kg}.$$
			\itemch 
			Gọi $x_1$, $x_2$, \ldots, $x_{42}$ là cân  nặng của $42$ học sinh được xếp theo thứ tự không giảm.\\
			Số trung vị của mẫu là $\dfrac{1}{2}\left(x_{21}+x_{22}\right)$, do $x_{21}$, $x_{22}\in[50{,}5;55{,}5)$ nên nhóm chứa trung vị là nhóm $[50{,}5;55{,}5)$.
			\itemch 
			Gọi $x_1$, $x_2$, \ldots, $x_{42}$ là cân  nặng của $42$ học sinh được xếp theo thứ tự không giảm.\\
			Tứ phân vị thứ nhất $Q_1$ là $x_{11}$ thuộc nhóm $[45{,}5;50{,}5)$.\\
			Ta có $Q_1=45{,}5 +\dfrac{\dfrac{42}{4} - 9}{7}\cdot (50{,}5 - 45{,}5)\approx 46{,}57$.
		\end{itemchoice}
	}
\end{ex}

\begin{ex}%[1D5H2-3]
% nguon: Archive/2024/Update-SP-BG10-11-3/SP-3/BG11-3/data_1C5/C5-B1-Dang4-5-PhanAnh.tex #1444
%[Bài Giảng 10-11 - Phan Anh]
	Thống kê cân nặng của $44$ học sinh lớp $11$A trong một trường THPT (đơn vị: kilôgam) được cho ở bẳng sau 
	\begin{center}
		\begin{tabular}{|c|c|c|}
			\hline
			Nhóm	&Tần số & Tần số tích lũy\\
			\hline
			$[30;40)$	&$4$ &$4$\\
			\hline
			$[40;50)$	&$10$ &$14$\\
			\hline
			$[50;60)$	&$18$ &$32$\\
			\hline
			$[60;70)$	&$10$ &$42$\\
			\hline
			$[70;80)$	&$2$ &$44$\\
			\hline
		\end{tabular}
	\end{center}
	Xét tính đúng - sai của các khẳng định sau
	\choiceTF[t]
	{Số học sinh cân nặng từ $40$ kg đến dưới $50$ kg là $18$}
	{\True $[50;60)$ là nhóm chứa số trung vị}
	{\True Cân nặng trung bình của các bạn học sinh trong lớp $11$A là $54$ kg}
	{Tứ phân vị thứ nhất của mẫu số liệu ghép nhóm trên là $50$ kg}
	\loigiai{
		\begin{itemchoice}
			\itemch 
			Số học sinh có cân nặng từ $40$ kg đến dưới $50$ kg là $10$.
			\itemch 
			Gọi $x_1$, $x_2$, \ldots, $x_{44}$ là cân nặng của $44$ học sinh được xếp theo thứ tự không giảm.\\
			Số trung vị của mẫu số liệu là $\dfrac{1}{2}\left(x_{22}+x_{23}\right)$, do $x_{22}$ và $x_{23}\in[50;60)$ nên nhóm $[50;60)$ chứa số trung vị của mẫu số liệu.
			\itemch 
			Cân nặng trung bình của học sinh lớp $11$A là
			$$\overline{x}=\dfrac{4\cdot 35+10\cdot 45+18\cdot 55+10\cdot 65+2\cdot 75}{44}\approx 54\,\text{(kg)}.$$
			\itemch 
			Số phần tử của mẫu là $n=44$.\\
			Ta có $\dfrac{n}{4}=\dfrac{44}{4}=11$. Mà $14>11$. Suy ra nhóm $2$ là nhóm đầu tiên có tần số tích lũy lớn hơn hoặc bằng $11$.\\
			Xét nhóm $2$ là nhóm $[40;50)$ có $s=40$;$h=10$;$n_2=10$ và nhóm $1$ là nhóm có $cf_1=4$.\\
			Ta có tứ phân vị thứ nhất là $Q_1=40 +\dfrac{11-4}{10}\cdot 10=47\,\text{(kg)}$.
		\end{itemchoice}
	}
\end{ex}

\begin{ex}%[1D5V2-3]
% nguon: Archive/2024/Update-SP-BG10-11-3/SP-3/BG11-3/data_1C5/C5-B1-Dang4-5-PhanAnh.tex #1445
%[Bài Giảng 10-11 - Phan Anh]
	Khảo sát cân nặng (đơn vị: kilôgam) của $40$ học sinh lớp 11D, thầy giáo chủ nhiệm thu được mẫu số liệu ghép nhóm sau
	\begin{center}
		\begin{tabular}{|c|c|c|c|c|c|c|}
			\hline 
			\begin{tabular}{c} 
				Cân năng\\
				(kilôgam)
			\end{tabular} & {$[30;40)$} & {$[40;50)$} & {$[50;60)$} & {$[60;70)$} & {$[70;80)$} &\\
			\hline Số hoc sinh & 3 & 18 & 12 & 5 & 2 & $n=40$\\
			\hline
		\end{tabular} 
	\end{center}
	\choiceTF[t]
	{\True Số lượng học sinh có cân nặng ít hơn $60$ kilôgam là $33$ học sinh}
	{Số lượng học sinh có cân nặng không dưới $50$ kilôgam là $12$ học sinh}
	{\True Trung vị của mẫu số liệu (làm tròn đến hàng phần chục) là $49{,}4$ kilôgam}
	{\True Trong số $40$ học sinh lớp 11D nói trên, có khoảng $25\%$ bạn trong lớp có cân nặng không dưới $57,5$ kilôgam và có khoảng $25\%$ số bạn trong lớp có cân nặng ít hơn $43{,}9$ kilôgam}
	\loigiai{
		\begin{itemchoice}
			\itemch 
			Số lượng học sinh có cân nặng ít hơn $60$ kilôgam là $3+18+12=33$ học sinh.
			\itemch 
			Số lượng học sinh có cân nặng không dưới $50$ kilôgam là $12+5+2=19$ học sinh.
			\itemch 
			Gọi $x_1, x_2,\ldots, x_{40}$ là cân nặng của $40$ học sinh lớp 11D xếp theo thứ tự không giảm.\\
			Số trung vị của mẫu số liệu trên là $\dfrac{1}{2}\left(x_{20}+x_{21}\right)$, nhận thấy $\dfrac{1}{2}\left(x_{20}+x_{21}\right)$ thuộc $[40;50)$.\\
			Áp dụng công thức với $n=40, n_m=18, C=3, u_m=40, u_{m+1}=50$, trung vị của mẫu số liệu ghép nhóm là
			$$
			M_e=u_m+\dfrac{\dfrac{n}{2}-C}{{{n}_{m}}}\left( u_{m+1}-u_m\right)=40+\frac{\dfrac{40}{2}-3}{18}\cdot(50-40)\approx49{,}4 .
			$$
			\itemch 
			Ta có tứ phân vị thứ nhất của dãy số liệu $x_1, x_2,\ldots, x_{40}$ là $\dfrac{1}{2}\left(x_{10}+x_{11}\right)x$. Vì $x_{10}, x_{11}\in[40;50)$ nên tứ phân vị thứ nhất của mẫu số liệu là 
			$$
			Q_1=u_i+\dfrac{\dfrac{n}{4}-C}{n_i}\left(u_{i+1}-u_i\right)=40+\dfrac{\dfrac{40}{4}-3}{18}(50-40)\approx 43{,}9.
			$$
			Do đó, có khoảng $25\%$ số bạn trong lớp có cân nặng ít hơn $43{,}9$ kilôgam.\\
			Tứ phân vị thứ ba của dãy số liệu $x_1, x_2,\ldots, x_{40}$ là $\dfrac{1}{2}\left(x_{30}+x_{31}\right)$. Vì $x_{30}, x_{31}\in[50;60)$ nên tứ phân vị thứ ba của mẫu số liệu là 
			$$
			Q_3=u_j+\frac{\dfrac{3 n}{4}-C}{n_j}\left(u_{j+1}-u_j\right)=50+\dfrac{\dfrac{3\cdot 40}{4}-(3+18)}{12}(60-50)=57{,}5.
			$$
			Do đó, có khoảng $25\%$ bạn trong lớp có cân nặng không dưới $57,5$ kilôgam.
		\end{itemchoice}
	}
\end{ex}

\begin{bt}%[1D5H2-3]
% nguon: Archive/2024/Update-SP-BG10-11-3/SP-3/BG11-3/data_1C5/C5-B1-Dang4-5-PhanAnh.tex #1447
[Mức độ 2]%[Bài Giảng 10-11 - Phan Anh]
	Doanh thu bán hàng trong $20$ ngày được lựa chọn ngẫu nhiên của một cửa hàng được ghi lại ở bảng sau (đơn vị: triệu đồng)
	\begin{center}
		\begin{tabular}{|c|c|c|c|c|c|}
			\hline Doanh thu & $[5;7)$ & $[7;9)$ & $[9;11)$ & $[11;13)$ & $[13;15)$\\
			\hline Số ngày  & $2$    & $7$    & $7$    & $3$     & $1$    \\
			\hline
		\end{tabular}
	\end{center}
	Tìm tứ phân vị thứ ba của mẫu số liệu trên, làm tròn đến hàng phần chục.
	\shortans{$10{,}7$}
	\loigiai{
		Tổng số ngày là $20$.\\
		Gọi $x_1;\ldots;x_{20}$ là doanh thu của cửa hàng trong $20$ ngày sắp xếp theo thứ tự không giảm.\\
		Ta có: $x_1;x_2\in[5;7), x_3;\ldots;x_9\in[7;9), x_{10};\ldots;x_{16}\in[9;11), x_{17};x_{18};x_{19}\in[11;13), x_{20}\in[13;15)$.\\
		Tứ phân vị thứ ba là $\dfrac{1}{2}\left(x_{15}+x_{16}\right)$ và $x_{15};x_{16}\in[9;11)$ nên ta có
		$$Q_2=9+\dfrac{\dfrac{3.20}{4}-9}{7}\cdot(11-9)\approx 10{,}7.$$}
\end{bt}

\begin{bt}%[1D5H2-3]
% nguon: Archive/2024/Update-SP-BG10-11-3/SP-3/BG11-3/data_1C5/C5-B1-Dang4-5-PhanAnh.tex #1449
[Mức độ 2]%[Bài Giảng 10-11 - Phan Anh]
	Phỏng vấn một số học sinh khối $11$ về thời gian ngủ (giờ) của một buổi tối thu được kết quả sau
	\begin{center}
		\begin{tabular}{|>{\centering\arraybackslash}p{3cm}*{5}{|>{\centering\arraybackslash}p{1.5cm}}|}
			\hline
			Thời gian & $[4;5)$ & $[5;6)$ & $[6;7)$ & $[7;8)$ & $[8;9)$\\
			\hline
			Số học sinh & $10$ & $18$ & $23$ & $20$ & $15$\\
			\hline
		\end{tabular}
	\end{center}
	Tìm thời gian ngủ ít nhất chiếm $75\%$ số học sinh ngủ trong một buổi tối, làm tròn đến hàng phần trăm.
	\shortans{$5{,}64$}
	\loigiai{
		Để tìm thời gian ngủ ít nhất của $75\%$ số học sinh trong một buổi tối, ta đi tìm phân vị thứ nhất của mẫu số liệu trên.\\
		Do cỡ mẫu là $n=86$ nên nhóm chứa giá trị $Q_1$ là $[5;6)$.\\
		Suy ra $Q_1=5+\dfrac{\dfrac{86}{4}-10}{18}\cdot 1\approx 5{,}64$.\\
		Vậy tứ phân vị thứ nhất của mẫu số liệu trên là $Q_1\approx 5{,}64$ nên $75\%$ số học sinh ngủ ít nhất $5{,}64$ giờ trong một buổi tối.
	}
\end{bt}

\begin{bt}%[1D5H1-4]
% nguon: Archive/2024/Update-SP-BG10-11-3/SP-3/BG11-3/data_1C5/C5-B1-Dang4-5-PhanAnh.tex #1450
[Mức độ 2]%[Bài Giảng 10-11 - Phan Anh]
	Kết quả kiểm tra môn Toán của lớp $11D$ như sau
	\[
	\begin{array}{cccccccccccccccccccc}
		5 & 6 & 7 & 5 & 6 & 9 & 10 & 8 & 5 & 5 & 4 & 5 & 4 & 5 & 7 & 4 & 5 & 8 & 9 & 10\\
		5 & 3 & 5 & 6 & 5 & 7 & 5 & 8 & 4 & 9 & 5 & 6 & 5 & 6 & 8 & 8 & 7 & 9 & 7 & 9
	\end{array}
	\]
	\begin{enumerate}
		\item Lập bảng tần số ghép nhóm của mẫu số liệu trên có bốn nhóm ứng với bốn nửa khoảng $\left[3;5\right)$, $\left[5;7\right)$, $\left[7;9\right)$, $\left[9;11\right)$.
		\item Mốt của bảng số liệu ghép nhóm trên là bao nhiêu (làm tròn kết quả đến hàng phần mười)?
	\end{enumerate}
	\loigiai{
		\immini
		{
			\begin{enumerate}
				\item Bảng bên là bảng tần số ghép nhóm cho kết quả kiểm tra môn Toán của lớp $11D$.
				\item Ta thấy: Nhóm $2$ ứng với nửa khoảng $\left[5;7\right)$ là nhóm có tần số lớn nhất với $u=5$, $g=2$, $n_2=18$. Nhóm $1$ có tần số $n_1=5$, nhóm $3$ có tần số $n_3=10$.\\
				Khi đó, mốt của mẫu số liệu là 
				\[
				M_o=5 +\left(\dfrac{18- 5}{2\cdot 18 - 5 - 10}\right)\cdot 2\approx 6{,}2.
				\]
			\end{enumerate}
		}
		{
			\begin{tabular}{|c|c|}
				\hline
				\textbf{Nhóm} &\textbf{Tần số}\\ 
				\hline
				$\left[3;5\right)$ & $5$\\
				\hline
				$\left[5;7\right)$ & $18$\\
				\hline
				$\left[7;9\right)$ & $10$\\
				\hline
				$\left[9;11\right)$ & $7$\\
				\hline
				& $n=40$\\
				\hline
			\end{tabular}
		}
	}
\end{bt}

\begin{ex}%[1D5H1-4]
% nguon: Archive/2024/Update-SP-BG10-11-3/SP-3/BG11-3/data_1C5/C5-B1-Dang4-5-PhanAnh.tex #1451
[Mức độ 2]%[Bài Giảng 10-11 - Phan Anh]
	Dựa vào điểm trung bình cuối năm học 2022-2023 của $487$ học sinh Khối $11$ đạt danh hiệu\,\lq\lq Học sinh Giỏi cả năm\rq\rq, ta thu được bảng số liệu sau
	\begin{center}
		\begin{tabular}{|c|c|c|c|c|c|c|c|c|}
			\hline
			\textbf{Điểm trung bình}& $[8{,}0;8{,}2)$ & $[8{,}2;8{,}4)$ & $[8{,}4;8{,}6)$ & $[8{,}6;8{,}8)$ & $[8{,}8;9{,}0)$ & $[9{,}0;9{,}2)$ & $[9{,}2;9{,}4)$ & $[9{,}4;9{,}6)$\\
			\hline
			\textbf{Số học sinh}	& $12$ & $31$ & $66$ & $107$ & $114$ & $93$ & $51$ & $13$ \\
			\hline
		\end{tabular}
	\end{center}
	Các mệnh đề sau đúng hay sai?
	\choiceTF[t]
	{\True Số trung bình của mẫu số liệu làm tròn đến hàng phần trăm là $8{,}84$}
	{\True Nhóm chứa mốt là $[8{,}8;9{,}0)$}
	{\True Trung vị của mẫu số liệu thuộc nhóm $[8{,}8;9{,}0)$}
	{Mốt của mẫu số liệu là $M_0=7{,}66$}
	\loigiai 
	{\begin{itemchoice}
			\itemch 
			Số trung bình của mẫu số liệu là
			\[\overline{x}=\dfrac{12\cdot8{,}1+31\cdot8{,}3+66\cdot8{,}5+107\cdot8{,}7+114\cdot8{,}9+93\cdot9{,}1+51\cdot9{,}3+13\cdot9{,}5}{12+31+66+107+114+93+51+13}=\dfrac{43\,051}{4\,870}\approx8{,}84.\]
			\itemch 
			Nhóm $[8{,}8;9{,}0)$ có tần số lớn nhất ($114$) nên nhóm chứa mốt của mẫu số liệu là $[8{,}8;9{,}0)$.
			\itemch 
			Gọi $x_1$, $x_2$, \ldots, $x_{487}$ là số điểm trung bình của $487$ học sinh xếp theo thứ tự không giảm.\\
			Trung vị của mẫu số liệu là $x_{244}\in[8{,}8;9{,}0)$.
			\itemch 
			Nhóm chứa mốt của mẫu số liệu là $[8{,}8;9{,}0)$\\
			Mốt của mẫu số liệu là
			\[M_0=8{,}8+\dfrac{114-107}{114-107+114-93}\cdot 0{,}2=8{,}85.\]
		\end{itemchoice}}
\end{ex}

\begin{bt}%[1D5H1-4]
% nguon: Archive/2024/Update-SP-BG10-11-3/SP-3/BG11-3/data_1C5/C5-B1-Dang4-5-PhanAnh.tex #1452
[Mức độ 2]%[Bài Giảng 10-11 - Phan Anh]
	Mẫu số liệu dưới đây ghi lại tốc độ của $40$ ô tô khi đi qua một trạm đo tốc độ (đơn vị: km/h)
	\[
	\begin{array}{cccccccccc}
		48{,}5 & 43 & 50 & 55 & 45 & 60 & 53 & 55,5 & 44 & 65\\
		51 & 62,5 & 41 & 44,5 & 57 & 57 & 68 & 49 & 46{,}5 & 53{,}5\\
		61 & 49{,}5 & 54 & 62 & 59 & 56 & 47 & 50 & 60 & 61\\
		49{,}5 & 52{,}5 & 57 & 47 & 60 & 55 & 45 & 47,5 & 48 & 61{,}5
	\end{array}
	\]
	\begin{enumerate}
		\item Lập bảng tần số ghép nhóm cho mẫu số liệu trên có sáu nhóm ứng với sáu nửa khoảng
		\[
		[40;45),[45;50),[50;55),[55;60),[60;65),[65;70).
		\]
		\item Mốt của mẫu số liệu ghép nhóm trên là bao nhiêu?
	\end{enumerate}
	\loigiai{
		\begin{enumerate}
			\item Ta có bảng tần số ghép nhóm của mẫu số liệu trên như sau
			\begin{center}
				\begin{tabular}{|c|c|c|c|}
					\hline
					\textbf{Nhóm} &\textbf{Giá trị đại diện} &\textbf{Tần số} &\textbf{Tần số tích luỹ}\\ 
					\hline
					$\left[40;45\right)$ & $42{,}5$ & $4$ & $4$\\
					$\left[45;50\right)$ & $47{,}5$ & $11$ & $15$\\
					$\left[50;55\right)$ & $52{,}5$ & $7$ & $22$\\
					$\left[55;60\right)$ & $57{,}5$ & $8$ & $30$\\
					$\left[60;65\right)$ & $62{,}5$ & $8$ & $38$\\
					$\left[65;70\right)$ & $67{,}5$ & $2$ & $40$\\
					\hline
					& & $n=40$ &\\
					\hline
				\end{tabular}
			\end{center}
			\item Ta thấy: Nhóm $2$ ứng với nửa khoảng $\left[45;50\right)$ là nhóm có tần số lớn nhất với $u=45$, $g=5$, $n_2=11$. Nhóm $1$ có tần số $n_1=4$, nhóm $3$ có tần số $n_3=7$.\\
			Khi đó, mốt của mẫu số liệu là 
			\[
			M_o=45 +\left(\dfrac{11 - 4}{2\cdot 11 - 4 - 7}\right)\cdot 5\approx 48{,}2\text{ (km/h)}.
			\]
		\end{enumerate}
	}
\end{bt}

\begin{bt}%[1D5H1-4]
% nguon: Archive/2024/Update-SP-BG10-11-3/SP-3/BG11-3/data_1C5/C5-B1-Dang4-5-PhanAnh.tex #1453
[Mức độ 2]%[Bài Giảng 10-11 - Phan Anh]
	Mẫu số liệu sau ghi lại cân nặng của $30$ bạn học sinh (đơn vị: kilôgam)
	\[
	\begin{array}{cccccccccc}
		17 & 40 & 39 & 40{,}5 & 42 & 51 & 41{,}5 & 39 & 41 & 30\\
		40 & 42 & 40{,}5 & 39{,}5 & 41 & 40{,}5 & 37 & 39{,}5 & 40 & 41\\
		38{,}5 & 39{,}5 & 40 & 41 & 39 & 40{,}5 & 40 & 38{,}5 & 39{,}5 & 41{,}5
	\end{array}
	\]
	\begin{enumerate}
		\item Lập bảng tần số ghép nhóm cho mẫu số liệu trên có tám nhóm ứng với tám nửa khoảng
		\[
		[15;20),[20;25),[25;30),[30;35),[35;40),[40;45),[45;50),[50;55).
		\]
		\item Mốt của mẫu số liệu ghép nhóm trên là bao nhiêu?
	\end{enumerate}
	\loigiai{
		\begin{enumerate}
			\item Ta có bảng tần số ghép nhóm của mẫu số liệu trên như sau
			\begin{center}
				\begin{tabular}{|c|c|c|c|}
					\hline
					\textbf{Nhóm} &\textbf{Giá trị đại diện} &\textbf{Tần số} &\textbf{Tần số tích luỹ}\\ 
					\hline
					$\left[15;20\right)$ & $17{,}5$ & $1$ & $1$\\
					$\left[20;25\right)$ & $22{,}5$ & $0$ & $1$\\
					$\left[25;30\right)$ & $27{,}5$ & $0$ & $1$\\
					$\left[30;35\right)$ & $32{,}5$ & $1$ & $2$\\
					$\left[35;40\right)$ & $37{,}5$ & $10$ & $12$\\
					$\left[40;45\right)$ & $42{,}5$ & $17$ & $29$\\
					$\left[45;50\right)$ & $47{,}5$ & $0$ & $29$\\
					$\left[50;55\right)$ & $52{,}5$ & $1$ & $30$\\
					\hline
					& & $n=30$ &\\
					\hline
				\end{tabular}
			\end{center}
			\item Ta thấy: Nhóm $6$ ứng với nửa khoảng $\left[40;45\right)$ là nhóm có tần số lớn nhất với $u=40$, $g=5$, $n_6=17$. Nhóm $5$ có tần số $n_5=10$, nhóm $7$ có tần số $n_7=0$.\\
			Khi đó, mốt của mẫu số liệu là 
			\[
			M_o=40 +\left(\dfrac{17 - 10}{2\cdot 17 - 10 - 0}\right)\cdot 5\approx 41{,}5\text{ (kg)}.
			\]
		\end{enumerate}
	}
\end{bt}

\begin{bt}%[1D5H1-4]
% nguon: Archive/2024/Update-SP-BG10-11-3/SP-3/BG11-3/data_1C5/C5-B1-Dang4-5-PhanAnh.tex #1454
[Mức độ 2]%[Bài Giảng 10-11 - Phan Anh]
	\immini
	{
		Bảng bên cho ta bảng tần số ghép nhóm số liệu thống kê chiều cao của $40$ mẫu cây ở một vườn thực vật (đơn vị: centimét).
		Mốt của mẫu số liệu ghép nhóm trên là bao nhiêu?
	}
	{
		\begin{tabular}{|c|c|c|}
			\hline
			\textbf{Nhóm} &\textbf{Tần số} &\textbf{Tần số tích luỹ}\\ 
			\hline
			$\left[30;40\right)$ & $4$ & $4$\\
			$\left[40;50\right)$ & $10$ & $14$\\
			$\left[50;60\right)$ & $14$ & $28$\\
			$\left[60;70\right)$ & $6$ & $34$\\
			$\left[70;80\right)$ & $4$ & $38$\\
			$\left[80;90\right)$ & $2$ & $40$\\
			\hline
			& $n=40$ &\\
			\hline
		\end{tabular}
	}
	\loigiai{
		Ta nhận thấy, nhóm $3$ ứng với nửa khoảng $\left[50;60\right)$ là nhóm có tần số lớn nhất với $u=50$, $g=10$, $n_3=14$. Nhóm $2$ có tần số $n_2=10$, nhóm $4$ có tần số $n_4=6$.\\
		Khi đó, mốt của mẫu số liệu là 
		\[
		M_o=50 +\left(\dfrac{14 - 10}{2\cdot 14 - 10 - 6}\right)\cdot 10\approx 53{,}3\text{ (cm)}.
		\]
	}
\end{bt}

\begin{ex}%[1D5H1-4]
% nguon: Archive/2024/Update-SP-BG10-11-3/SP-3/BG11-3/data_1C5/C5-B1-Dang4-5-PhanAnh.tex #1457
[Mức độ 2]%[Bài Giảng 10-11 - Phan Anh]
	Thâm niên giảng dạy của một số giáo viên trường THPT được ghi lại ở bảng sau
	\begin{center}
		\begin{tabular}{|c|c|c|c|c|c|}
			\hline Thâm niên (Số năm) &$[1; 5)$ &$[5; 10)$ &$[10; 15)$ &$[15; 20)$ &$[20; 25)$ \\
			\hline Số giáo viên & 4 & 12 & 16 & 8 & 3 \\
			\hline
		\end{tabular}
	\end{center}
	\choiceTF[t]
	{\True Cỡ mẫu của mẫu số liệu bằng $50$}
	{\True Số trung bình của mẫu ghép nhóm làm tròn đến hàng phần trăm là $11{,}85$}
	{\True Nhóm chứa mốt của mẫu số liệu trên là nhóm $[10;15)$}
	{Mốt của mẫu số liệu ghép nhóm làm tròn đến hàng phần trăm là $11{,}74$}
	\loigiai{Ta có bảng số liệu như sau
		\begin{center}
			\begin{tabular}{|c|c|c|c|c|c|}
				\hline Thâm niên (Số năm) &$[1; 5)$ &$[5; 10)$ &$[10; 15)$ &$[15; 20)$ &$[20; 25)$ \\
				\hline Giá trị đại diện & 3 & 7{,}5 & 12{,}5 & 17{,}5 & 22{,}5 \\
				\hline Số giáo viên & 4 & 12 & 16 & 8 & 3 \\
				\hline
			\end{tabular}
		\end{center}
		\begin{itemchoice}
			\itemch 
			Cỡ mẫu của mẫu số liệu bằng $n=4+12+16+8+3=43$.
			\itemch 
			Số trung bình của mẫu ghép là
			\[\overline{x}=\dfrac{3\cdot 4+7{,}5\cdot 12+12{,}5\cdot 16+17{,}5\cdot 8+22{,}5\cdot 3}{4+12+16+8+3}=\dfrac{1\,019}{86}\approx11{,}85.\]
			\itemch 
			Nhóm chứa mốt của mẫu số liệu trên là nhóm $[10;15)$.
			\itemch 
			Ta có $u_m=10$, $n_{m-1}=12$, $n_m=16$, $n_{m+1}=8$, $n_{m+1}-u_m=15-10=5$.
			\begin{center}
				$M_{0}=10+\dfrac{16-12}{(16-12)+(16-8)} \cdot 5=\dfrac{35}{3}\approx11{,}67$.
			\end{center}
	\end{itemchoice}}
\end{ex}

\begin{ex}%[1D5H1-4]
% nguon: Archive/2024/Update-SP-BG10-11-3/SP-3/BG11-3/data_1C5/C5-B1-Dang4-5-PhanAnh.tex #1458
[Mức độ 2]%[Bài Giảng 10-11 - Phan Anh]
	Một nhà thực vật học đo chiều dài trung bình của $74$ lá cây (đơn vị milimét) và thu được bảng tần số ghép nhóm như sau
	\begin{center}
		\begin{tabular}{|c|c|c|}
			\hline Nhóm & Giá trị đại diện & Tần số\\
			\hline$[5{,}45;5{,}85)$ & $5{,}65$ & $5$\\
			\hline$[5{,}85;6{,}25)$ & $6{,}05$ & $9$\\
			\hline$[6{,}25;6{,}65)$ & $6{,}45$ & $15$\\
			\hline$[6{,}65;7{,}05)$ & $6{,}85$ & $19$\\
			\hline$[7{,}05;7{,}45)$ & $7{,}25$ & $16$\\
			\hline$[7{,}45;7{,}85)$ & $7{,}65$ & $8$\\
			\hline$[7{,}85;8{,}25)$ & $8{,}05$ & $2$\\
			\hline
		\end{tabular}
	\end{center}
	\choiceTF[t]
	{Chiều dài trung bình của $74$ lá cây bằng $\approx 6{,}4$ (mm)}
	{\True Độ dài nhóm là $0{,}4$} 
	{Nhóm chứa mốt là $[7{,}05;7{,}45)$} 
	{mốt của mẫu số liệu ghép nhóm là $\approx 6{,}65$} 
	\loigiai { 
		\begin{itemchoice}
			\itemch	
			Chiều dài trung bình của $74$ lá cây là 		
			\begin{center}
				$\overline{x}=\dfrac{5{,}65\cdot5+6{,}05\cdot9+6{,}45\cdot 15+6{,}85\cdot 19+7{,}25\cdot16+7{,}65\cdot 8+8{,}05\cdot2}{74}=\dfrac{5029}{740}\approx 6{,}8\,\text{(mm)}$.
			\end{center}
			\itemch	
			Độ dài nhóm là $0{,}4$.
			\itemch 
			Nhóm chứa mốt là $[6{,}65;7{,}05)$, trong đó
			\begin{center}
				$u_m=6{,}65;u_{m+1}=7{,}05;u_{m+1}-u_m=0{,}4;n_m=19;n_{m-1}=15;n_{m+1}=16$.
			\end{center}
			\itemch	
			Do đó mốt của mẫu số liệu ghép nhóm là
			$M_o=6{,}65+\dfrac{19-15}{(19-15)+(19-16)}\cdot 0{,}4\approx 6{,}88$.
		\end{itemchoice}
	}
\end{ex}

\begin{bt}%[1D5H1-4]
% nguon: Archive/2024/Update-SP-BG10-11-3/SP-3/BG11-3/data_1C5/C5-B1-Dang4-5-PhanAnh.tex #1459
[Mức độ 2]%[Bài Giảng 10-11 - Phan Anh]
	Khảo sát chiều cao của một số học sinh khối $11$ thu được mẫu số liệu ghép nhóm sau
	\begin{center}	
		\begin{tabular}{|l|c|c|c|c|c|}
			\hline
			Khoảng chiều cao (cm)&	$[145;150)$& $[150;155)$&$[155;160)$&$[160;165)$& $[165;170)$\\
			\hline
			Số học sinh&	$7$&$14$&$10$&$10$ &$9$\\\hline
		\end{tabular}
	\end{center} 
	Tính mốt của mẫu số liệu ghép nhóm này (làm tròn kết quả đến hàng đơn vị).
	\shortans{$153$}
	\loigiai{
		Tần số lớn nhất là $14$ nên nhóm chứa mốt là nhóm $[150;155)$.\\
		Ta có nhóm có tần số lớn nhất là nhóm $i=2$;giá trị bên trái của nhóm $2$ là $a_2=150$ với tần số $n_2=14$;tần số nhóm trước nó là $n_1=7$ và tần số nhóm sau là $n_3=10$;độ dài nhóm $2$ là $h=5$.\\
		Do đó $M_0=a_2+\left(\dfrac{n_i-n_{i-1}}{2n_i-n_{i-1}-n_{i+1}}\cdot h\right) =150+\dfrac{14-7}{(14-7)+(14-10)}\cdot5=\dfrac{1685}{11}\approx 153$.
	}
\end{bt}

\begin{bt}%[1D5N1-4]
% nguon: Archive/2024/Update-SP-BG10-11-3/SP-3/BG11-3/data_1C5/C5-B1-Dang4-5-PhanAnh.tex #1460
[Mức độ 2]%[Bài Giảng 10-11 - Phan Anh]
	Khảo sát thời gian tập thể dục trong ngày của một số học sinh khối 11 thu được mẫu số liệu ghép nhóm sau
	\begin{center}
		\begin{tabular}{|c|c|c|c|c|c|}
			\hline
			\textbf{Thời gian (phút)}&$[0; 20)$&$[20; 40)$&$[40; 60)$&$[60; 80)$&$[80; 100)$\\
			\hline
			\textbf{Số học sinh}&$5$&$9$&$12$&$10$&$6$\\
			\hline
		\end{tabular}
	\end{center}
	Tìm mốt của mẫu số liệu ghép nhóm trên.
	\shortans{$52$}
	\loigiai{
		Nhóm $[40; 60)$ có tần số bằng $12$ là lớn nhất nên nhóm chứa mốt của mẫu số liệu là nhóm $[40; 60)$.
		Do đó $u_m=40$, $u_{m+1}=60$, $n_m=12$, $n_{m-1}=9$, $n_{m+1}=10$\\
		Vậy mốt của mẫu số số liệu là
		\[M_0=u_m+\dfrac{n_m-n_{m-1}}{\left(n_m-n_{m-1}\right)+\left(n_m-n_{m+1}\right)}\cdot\left(u_{m+1}-u_m\right)=40+\dfrac{12-9}{(12-9)+(12-10)}\cdot (60-40)=52.\]
	}
\end{bt}

\begin{bt}%[1D5H1-4]
% nguon: Archive/2024/Update-SP-BG10-11-3/SP-3/BG11-3/data_1C5/C5-B1-Dang4-5-PhanAnh.tex #1461
[Mức độ 2]%[Bài Giảng 10-11 - Phan Anh]
	Doanh thu bán hàng trong 20 ngày được lựa chọn ngẫu nhiên của một cửa hàng được ghi lại ở bảng sau (đơn vị: triệu đồng)
	\begin{center}
		\begin{tabular}{|c|c|c|c|c|c|}
			\hline
			Doanh thu& $[5;7)$ & $[7;9)$ & $[9;11)$ & $[11;13)$ & $[13;15)$\\
			\hline
			Số ngày&$2$ &$7$ &$7$ &$3$ &$1$\\
			\hline
		\end{tabular}
	\end{center}
	Tìm mốt của mẫu số liệu ghép nhóm trên.
	\shortans{$9$}
	\loigiai{
		Nhóm chứa mốt của mẫu số liệu trên là nhóm $[7;9)$ hoặc $[9;11)$.
		\begin{enumerate}
			\item [{\bf TH1.}] Xét nhóm $[7;9)$ ta có $u_m=7$, $u_{m+1}=9$, $n_m=7$, $n_{m+1}=7$, $n_{m-1}=2$.\\
			Mốt của mẫu số liệu ghép nhóm là 
			\allowdisplaybreaks{
				\begin{eqnarray*}
					M_o&=&u_m+\dfrac{n_m-n_{m-1}}{\left(n_m-n_{m-1}\right)+\left(n_m-n_{m+1}\right)}\cdot\left(u_{m+1}-u_m\right)\\
					&=& 7+\dfrac{7-2}{\left(7-2\right)+\left(7-7\right)}\cdot\left(9-7\right)=9.
			\end{eqnarray*} }
			\item [{\bf TH2.}] Xét nhóm $[9;11)$ ta có $u_m=9$, $u_{m+1}=11$, $n_m=7$, $n_{m+1}=3$, $n_{m-1}=7$.\\
			Mốt của mẫu số liệu ghép nhóm là 
			\allowdisplaybreaks{
				\begin{eqnarray*}
					M_o&=&u_m+\dfrac{n_m-n_{m-1}}{\left(n_m-n_{m-1}\right)+\left(n_m-n_{m+1}\right)}\cdot\left(u_{m+1}-u_m\right)\\
					&=& 9+\dfrac{7-7}{\left(7-7\right)+\left(7-3\right)}\cdot\left(11-9\right)=9.
			\end{eqnarray*} }
		\end{enumerate}
		Vậy mốt của mẫu số liệu ghép nhóm trên là $M_o=9$.
	}
\end{bt}

\begin{bt}%[1D5H1-4]
% nguon: Archive/2024/Update-SP-BG10-11-3/SP-3/BG11-3/data_1C5/C5-B1-Dang4-5-PhanAnh.tex #1462
[Mức độ 2]%[Bài Giảng 10-11 - Phan Anh]
	Tìm hiểu thời gian hoàn thành một bài tập (đơn vị: phút) của một số học sinh thu được kết quả sau
	\begin{center}
		\begin{tabular}{|c|c|c|c|c|c|}
			\hline
			\begin{tabular}{l}
				Thời gian(giờ)
			\end{tabular}&{$[0;4)$}&{$[4;8)$}&{$[8;12)$}&{$[12;16)$}&{$[16;20)$}\\
			\hline Số học sinh & $2$ & $4$ & $7$ & $4$ & $3$\\
			\hline
		\end{tabular}
	\end{center}
	Tìm mốt của mẫu số liệu ghép nhóm trên.
	\shortans{$10$}
	\loigiai{Nhóm chứa Mốt của bảng số liệu này là $[8;12)$, suy ra Mốt của bảng số liệu là
		$$M_o=8+\dfrac{7-4}{(7-4)+(7-4)}\cdot(12-8)=10.$$}
\end{bt}

\begin{bt}%[1D5H1-4]
% nguon: Archive/2024/Update-SP-BG10-11-3/SP-3/BG11-3/data_1C5/C5-B1-Dang4-5-PhanAnh.tex #1463
[Mức độ 2]%[Bài Giảng 10-11 - Phan Anh]
	Các bạn học sinh lớp $ 11A1$ trả lời $ 40$ câu hỏi trong một bài kiểm tra. Kết quả được thống kê trong bảng tần số ghép nhóm sau
	\begin{center}
		\begin{tabular}{|c|c|c|c|c|c|}
			\hline
			Số câu trả lời đúng  & $\left[16;21\right)$ & $\left[21;26\right)$ & $\left[26;31\right)$ & $\left[31;36\right)$ & $\left[36;41\right)$\\
			\hline
			Số học sinh & $4$ & $6$ & $8$ & $18$ & $4$\\
			\hline
		\end{tabular}
	\end{center}
	Tìm mốt của mẫu số liệu ghép nhóm trên, làm tròn đến hàng phần chục.
	\shortans{$33{,}1$}
	\loigiai{
		Tần số lớn nhất là $18$ nên nhóm chứa mốt là nhóm $[31;36)$.\\
		Ta có $u_m=31, u_{m+1}=36, n_m=18, n_{m-1}=8, n_{m+1}=4$.\\
		Do đó $M_o=31+\dfrac{18-8}{(18-8)+(18-4)}\cdot (36-31)=\dfrac{397}{12}\approx 33{,}1$.
	}
\end{bt}

\begin{ex}%[1D5H1-3]
% nguon: Archive/2024/Update-SP-BG10-11-3/SP-3/BG11-3/data_1C5/ChuongV-Bai1-Dang2-SoTrungBinhCong-Diamon.tex #1510
 
Tuổi các học viên của một lớp học Tiếng Anh tại một trung tâm được ghi lại ở bảng tần số ghép lớp như sau
\begin{center}
\begin{tabular}{|w{c}{2cm}|w{c}{2cm}|w{c}{2cm}|w{c}{2cm}|w{c}{2cm}|w{c}{2cm}|}
\hline
Lớp & $[16; 20)$ & $[20; 24)$ & $[24; 28)$ & $[28; 32)$ & $[32; 36)$ \\
\hline
Tần số & $10$ & $12$ & $14$ & $9$ & $5$\\
\hline
\end{tabular}
\end{center}
Các khẳng định sau đúng hay sai?
\choiceTF[t]
{\True Tổng số học viên $50$}
{\True Độ dài của mỗi nhóm là $4$}
{Giá trị đại diện của nhóm $[24; 28)$ là $4$}
{\True Độ tuổi trung bình của học viên trung tâm là  $25$}
\loigiai{
Ta có bảng sau
\begin{center}
\begin{tabular}{|w{c}{2.7cm}|w{c}{2cm}|w{c}{2cm}|w{c}{2cm}|w{c}{2cm}|w{c}{2cm}|}
\hline
Lớp & $[16; 20)$ & $[20; 24)$ & $[24; 28)$ & $[28; 32)$ & $[32; 36)$ \\
\hline
Tần số & $10$ & $12$ & $14$ & $9$ & $5$\\
\hline
Giá trị đại diện&$18$&$22$&$26$&$30$&$34$\\
\hline
\end{tabular}
\end{center}
\begin{itemchoice}
\itemch Tổng số học viên là $10+12+14+9+5=50$.
\itemch Độ dài của nhóm là $20-16=4$.
\itemch  Giá trị đại diện của  nhóm $[24; 28)$ là $\dfrac{24+28}{2}=26$.
\itemch Độ tuổi trung bình của học viên trung tâm bằng
\[\overline{x} = \dfrac{10 \cdot 18 + 12 \cdot 22 + 14 \cdot 26 + 9 \cdot 30 + 5 \cdot 35}{50} \approx 25.\]
\end{itemchoice}
}
\end{ex}

\begin{ex}%[1D5H1-3]
% nguon: Archive/2024/Update-SP-BG10-11-3/SP-3/BG11-3/data_1C5/ChuongV-Bai1-Dang2-SoTrungBinhCong-Diamon.tex #1511

Khảo sát thời gian tập thể dục của một số học sinh khối 11 thu được mẫu số liệu ghép nhóm sau
\begin{center}
\begin{tabular}{|c|c|c|c|c|c|}
\hline
Thời gian (phút)&$[0;20)$  &$[20;40)$  &$[40;60)$  &$[6;80)$  &$[80;100)$  \\
\hline
Số học sinh&$5$  &$9$  &$12$  &$10$  &$6$  \\
\hline
\end{tabular}
\end{center}
Xét tính đúng - sai của các khẳng định sau:
\choiceTF[t]
{\True Tổng số học sinh được khảo sát là $42$ học sinh}
{Giá trị đại diện của nhóm $[20;40)$ là $25$}
{\True Có $16$ học sinh tập thể dục ít nhất $1$ giờ trong ngày}
{Số trung bình của mẫu số liệu trên thuộc nhóm $[20;40)$}
\loigiai{
\begin{itemchoice}
\itemch Tổng số học sinh được khảo sát là $n=5+9+12+10+6=42$.
\itemch  Giá trị đại diện của nhóm $[20;40)$ là $\dfrac{20+40}{2}=30$.
\itemch Số học sinh tập thể dục ít nhất $1$ giờ trong ngày là $10+6=16$ học sinh.
\itemch  Ta có
\begin{center}
\begin{tabular}{|c|c|c|c|c|c|}
\hline
Thời gian (phút)&$[0;20)$  &$[20;40)$  &$[40;60)$  &$[6;80)$  &$[80;100)$  \\
\hline
Giá trị đại diện &10&30&50&70&90\\
\hline
Tần số&$5$  &$9$  &$12$  &$10$  &$6$  \\
\hline
\end{tabular}
\end{center}
Cỡ mẫu là $n=5+9+12+10+6=42$.\\
Thời gian trung bình học sinh tập thể dục là
$$\overline{x} = \dfrac{5\cdot 10 + 9 \cdot 30 + 12 \cdot 50 + 10\cdot 70 + 6\cdot 90}{42} = \dfrac{360}{7} \in [40;60).$$
\end{itemchoice}
}
\end{ex}

\begin{ex}%[1D5H1-3]
% nguon: Archive/2024/Update-SP-BG10-11-3/SP-3/BG11-3/data_1C5/ChuongV-Bai1-Dang2-SoTrungBinhCong-Diamon.tex #1512
%[Dự án True-False-2024]%[Phan Văn Thành]
Thời gian tự học mỗi ngày của học sinh khối $11$ được thống kê trong bảng sau
\begin{center}
\begin{tabular}{|c|c|c|c|c|c|}
\hline
Thời gian (phút)&$[60;90)$  &$[90;120)$  &$[120;150)$  &$[150;180)$  &$[180;210)$  \\
\hline
Số học sinh&$8$  &$10$  &$13$  &$7$  &$2$  \\
\hline
\end{tabular}
\end{center}
Xét tính đúng - sai của các khẳng định sau:
\choiceTF[t]
{\True Tổng số học sinh được khảo sát là $40$}
{\True Giá trị đại diện của nhóm $[90;120)$ là $105$}
{Thời gian trung bình tự học mỗi ngày của các bạn học sinh được khảo sát là $135$ phút}
{Có $13$ học sinh tự học ít nhất $2$ giờ mỗi ngày}
\loigiai{
\begin{itemchoice}
\itemch Tổng số học sinh được khảo sát là $n=8+10+13+7+2 = 40$ (học sinh).
\itemch Giá trị đại diện của nhóm $[90;120)$ là $\dfrac{90+120}{2} = 105$.
\itemch  Ta có
\begin{center}
\begin{tabular}{|c|c|c|c|c|c|}
\hline
Thời gian (phút)&$[60;90)$  &$[90;120)$  &$[120;150)$  &$[150;180)$  &$[180;210)$  \\
\hline
Giá trị đại diện &75&105&135&165&195\\
\hline
Số học sinh&$8$  &$10$  &$13$  &$7$  &$2$  \\
\hline
\end{tabular}
\end{center}
Thời gian trung bình học sinh tự học là
$$\overline{x} = \dfrac{75\cdot 8 + 105 \cdot 10 + 135 \cdot 13 + 165\cdot 7 + 195\cdot 2}{40} = \dfrac{495}{4}\approx 123{,}75\, \text{(phút)}.$$
\itemch  Số học sinh tự học ít nhất $2$ giờ mỗi ngày là $13+7+2=22$ (học sinh).
\end{itemchoice}
}
\end{ex}

\begin{ex}%[1D5H1-4]
% nguon: Archive/2024/Update-SP-BG10-11-3/SP-3/BG11-3/data_1C5/ChuongV-Bai1-Dang2-SoTrungBinhCong-Diamon.tex #1513
%[Dự án True-False-2024]%[Nguyen Huynh]
Kết quả khảo sát cân nặng của $25$ quả cam ở mỗi lô hàng $A, B$ được cho ở bảng sau
\begin{center}
\begin{tabular}{|c|c|c|c|c|c|}
\hline Cân nặng (gam) & {$[150;155)$} & {$[155;160)$} & {$[160;165)$} & {$[165;170)$} & {$[170;175)$} \\
\hline Số quả cam ở lô hàng $A$ & $2$ & $6$ & $12$ & $4$ & $1$ \\
\hline Số quả cam ở lô hàng $B$ & $1$ & $3$ & $7$ & $10$ & $4$ \\
\hline
\end{tabular}
\end{center}
Xét tính đúng - sai của các khẳng định sau:
\choiceTF[t]
{\True Giá trị đại diện nhóm $[150;155)$ bằng $152{,}5$} 
{Cân nặng trung bình của mỗi quả cam ở lô $A$ là $163{,}7$ (gam)} 
{Cân nặng trung bình của mỗi quả cam ở lô $B$ là $162{,}1$ (gam)} 
{\True Theo số trung bình thì cam ở lô hàng $B$ nặng hơn cam ở lô hàng $A$} 
\loigiai{
Bảng thống kê số lượng cam theo giá trị đại diện
\begin{center}
\begin{tabular}{|c|c|c|c|c|c|}
\hline Cân nặng đại diện (gam) & $152{,}5$ & $157{,}5$ & $162{,}5$ & $167{,}5$ & $172{,}5$\\
\hline Số quả cam ở lô hàng $A$ & $2$ & $6$ & $12$ & $4$ & $1$\\
\hline Số quả cam ở lô hàng $B$ & $1$ & $3$ & $7$ & $10$ & $4$\\
\hline
\end{tabular}
\end{center}
\begin{itemchoice}
\itemch Giá trị đại diện nhóm $[150;155)$ bằng $152{,}5$.
\itemch  Cân nặng trung bình của mỗi quả cam ở lô $A$ là
\begin{center}
$\overline{x}_A=\dfrac{152{,}5 \cdot 2+157{,}5 \cdot 6+162{,}5 \cdot 12+167{,}5 \cdot 4+172{,}5 \cdot 1}{25}=161{,}7$.
\end{center}
\itemch  Cân nặng trung bình của mỗi quả cam ở lô $B$ là
\begin{center}
$\overline{x}_B=\dfrac{152{,}5 \cdot 1+157{,}5 \cdot 3+162{,}5 \cdot 7+167{,}5 \cdot 10+172{,}5 \cdot 4}{25}=165{,}1$.
\end{center}
\itemch Ta thấy $\overline{x}_A<\overline{x}_B$. Vậy nếu so sánh theo số trung bình thì cam ở lô hàng $B$ nặng hơn cam ở lô hàng $A$.
\end{itemchoice}
}
\end{ex}

\begin{bt}%[1D5H1-3]
% nguon: Archive/2024/Update-SP-BG10-11-3/SP-3/BG11-3/data_1C5/ChuongV-Bai1-Dang2-SoTrungBinhCong-Diamon.tex #1514
 
Cân nặng (đơn vị: kg) của học sinh lớp $11A$ được cho trong bảng dưới đây
\begin{center}
\begin{tabular}{|c|c|c|c|c|c|}
\hline
Cân nặng (kg)& $\left[40;45\right)$ & $\left[45;50\right)$  & $\left[50;55\right)$&$\left[55;60\right)$ &$\left[60;65\right)$    \\
\hline
Số học sinh& $8$ & $11$ &$9$  & $8$ & $4$ \\
\hline
\end{tabular}
\end{center}
Tính cân nặng trung bình của học sinh lớp $11A$ (làm tròn kết quả đến chữ số thập phân thứ nhất).
\shortans[]{$51{,}1$}
\loigiai{
\begin{center}
\renewcommand{\arraystretch}{1.4} 
\begin{tabular}{|c|c|c|c|c|c|}
\hline
Cân nặng (kg)& $\left[40;45\right)$ & $\left[45;50\right)$  & $\left[50;55\right)$&$\left[55;60\right)$ &$\left[60;65\right)$  \\
\hline
Giá trị đại diện cân nặng &$42{,}5$&$47{,}5$ &$52{,}5$ &$57{,}5$ &$62{,}5$
\\
\hline
Số học sinh& $8$ & $11$ &$9$  & $8$ & $4$ \\
\hline
\end{tabular}
\end{center}
Cân nặng trung bình của của học sinh lớp $11$A là
$$\overline{x}=\dfrac{42{,}5\cdot8+47{,}5\cdot11+52{,}5\cdot9+57{,}5\cdot8+62{,}5\cdot4}{8+11+9+8+4}\approx 51{,}1\; (\text{kg}).$$
}
\end{bt}

\begin{bt}%[1D5H1-3]
% nguon: Archive/2024/Update-SP-BG10-11-3/SP-3/BG11-3/data_1C5/ChuongV-Bai1-Dang2-SoTrungBinhCong-Diamon.tex #1515

Mức thưởng tết (triệu đồng) của các công nhân trong một nhà máy được cho trong bảng sau
\begin{center}
\begin{tabular}{|c|c|c|c|c|c|}
\hline 
\textbf{Mức thưởng} & $[5; 10)$ & $[10; 15)$ & $[15; 20)$ & $[20; 25)$ & $[25; 30)$ \\ 
\hline 
\textbf{Số công nhân} & $19$ & $30$ & $42$ & $6$ & $3$ \\ 
\hline 
\end{tabular}  
\end{center} 
Giá trị trung bình của mẫu số liệu ghép nhóm là (làm tròn đến hàng phần chục)
\shortans[]{$17{,}7$}
\loigiai{
Ta có
\begin{center}
\begin{tabular}{|c|c|c|c|c|c|}
\hline 
\textbf{Mức thưởng} & $[5; 10)$ & $[10; 15)$ & $[15; 20)$ & $[20; 25)$ & $[25; 30)$ \\ 
\hline
\textbf{Giá trị đại diện}&$7{,}5$&$12{,}5$&$17{,}5$&$22{,}5$&$27{,}5$\\
\hline
\textbf{Số công nhân} & $19$ & $30$ & $42$ & $6$ & $3$ \\ 
\hline 
\end{tabular}  
\end{center} 
Giá trị trung bình của mẫu số liệu là $$\overline{x}=\dfrac{19\cdot 7{,}5+30\cdot 12{,}5+42\cdot 17{,}5+6\cdot 22{,}5+3\cdot 27{,}5}{19+30+42+6+3}=14{,}7$$
}
\end{bt}

\begin{bt}%[1D5N1-3]
% nguon: Archive/2024/Update-SP-BG10-11-3/SP-3/BG11-3/data_1C5/ChuongV-Bai1-Dang2-SoTrungBinhCong-Diamon.tex #1516

Tuổi các học viên của một lớp học Tiếng Anh tại một trung tâm được ghi lại ở bảng tần số ghép lớp như sau
\begin{center}
\begin{tabular}{|c|c|c|c|c|c|}
\hline
Lớp & $[16;20)$ & $[20;24)$ & $[24;28)$ & $[28;32)$ & $[32;36)$ \\
\hline
Tần số & $10$ & $12$ & $14$ & $9$ & $5$\\ 
\hline
\end{tabular}
\end{center}
Độ tuổi trung bình của học viên trung tâm bằng bao nhiêu (làm tròn đến hàng phần chục)?
\shortans[]{$25$}
\loigiai{
Tổng số học viên $50$.
\begin{center}
\begin{tabular}{|c|c|c|c|c|c|}
\hline
Lớp & $[16;20)$ & $[20;24)$ & $[24;28)$ & $[28;32)$ & $[32;36)$ \\
\hline
Giá trị đại diện&$18$&$22$&$26$&$30$&$34$\\
\hline
Tần số & $10$ & $12$ & $14$ & $9$ & $5$\\
\hline
\end{tabular}
\end{center}
Độ tuổi trung bình của học viên trung tâm bằng
$$\overline{x}=\dfrac{10\cdot 18+22\cdot 12+14\cdot 26+9\cdot 30+5\cdot 34}{50}\approx 25.$$
}
\end{bt}

\begin{bt}%[1D5N1-3]
% nguon: Archive/2024/Update-SP-BG10-11-3/SP-3/BG11-3/data_1C5/ChuongV-Bai1-Dang2-SoTrungBinhCong-Diamon.tex #1517

Kết quả kiểm tra $1$ tiết môn Toán của một lớp học được ghi lại ở bảng tần số ghép lớp như sau
\begin{center}
\begin{tabular}{|c|c|c|c|c|c|}
\hline
Điểm & $[5;6)$ & $[6;7)$ & $[7;8)$ & $[8;9)$ & $[9;10]$\\
\hline
Tần số & $3$ & $10$ & $15$ & $8$ & $4$\\
\hline
\end{tabular}
\end{center}
Điểm trung bình môn Toán của lớp học bằng bao nhiêu (làm tròn đến hàng phần mười)?
\shortans[]{$7{,}5$}
\loigiai{
Tổng số học viên $n=3+10+15+8+4=40$.
\begin{center}
\begin{tabular}{|c|c|c|c|c|c|}
\hline
Điểm & $[5;6)$ & $[6;7)$ & $[7;8)$ & $[8;9)$ & $[9;10]$ \\
\hline
Giá trị đại diện&$5{,}5$&$6{,}5$&$7{,}5$&$8{,}5$&$9{,}5$\\
\hline
Tần số & $3$ & $10$ & $15$ & $8$ & $4$\\
\hline
\end{tabular}
\end{center}
Độ tuổi trung bình của học viên trung tâm bằng
$$\overline{x}=\dfrac{5{,}5\cdot 3+6{,}5\cdot 10+7{,}5\cdot 15+8{,}5\cdot 8+9{,}5\cdot 4}{40}= 7{,}5.$$
}
\end{bt}

\begin{bt}%[1D5N1-3]
% nguon: Archive/2024/Update-SP-BG10-11-3/SP-3/BG11-3/data_1C5/ChuongV-Bai1-Dang2-SoTrungBinhCong-Diamon.tex #1518

Thời gian (phút) mà một nhóm học sinh dành để hoàn thành một bài tập về nhà được ghi lại ở bảng tần số ghép lớp như sau
\begin{center}
\begin{tabular}{|c|c|c|c|c|c|}
\hline
Thời gian (phút) & $[20;30)$ & $[30;40)$ & $[40;50)$ & $[50;60)$ & $[60;70)$\\
\hline
Tần số & $5$ & $12$ & $18$ & $7$ & $3$\\
\hline
\end{tabular}
\end{center}
Thời gian trung bình mà một học sinh hoàn thành bài tập là bao nhiêu phút (làm tròn đến hàng đơn vị)?
\shortans[]{$43$}
\loigiai{
Tổng số học viên $n=5+12+18+7+3=45$.
\begin{center}
\begin{tabular}{|c|c|c|c|c|c|}
\hline
Thời gian (phút) & $[20;30)$ & $[30;40)$ & $[40;50)$ & $[50;60)$ & $[60;70)$\\
\hline
Giá trị đại diện&$25$&$35$&$45$&$55$&$65$\\
\hline
Tần số & $5$ & $12$ & $18$ & $7$ & $3$\\
\hline
\end{tabular}
\end{center}
Độ tuổi trung bình của học viên trung tâm bằng
$$\overline{x}=\dfrac{25\cdot 5+35\cdot 12+45\cdot 18+55\cdot 7+65\cdot 3}{45}= 43.$$
}
\end{bt}

\begin{bt}%[1D5N1-3]
% nguon: Archive/2024/Update-SP-BG10-11-3/SP-3/BG11-3/data_1C5/ChuongV-Bai1-Dang2-SoTrungBinhCong-Diamon.tex #1519

Chiều cao (cm) của các cây hoa trong một vườn được ghi lại ở bảng tần số ghép lớp như sau
\begin{center}
\begin{tabular}{|c|c|c|c|c|c|}
\hline
Chiều cao (cm) & $[10;20)$ & $[20;30)$ & $[30;40)$ & $[40;50)$ & $[50;60)$\\
\hline
Tần số & $2$ & $8$ & $15$ & $10$ & $5$\\
\hline
\end{tabular}
\end{center}
Chiều cao trung bình của các cây hoa trong vườn là bao nhiêu xentimet (làm tròn đến hàng đơn vị)?
\shortans[]{$37$}
\loigiai{
Tổng số học viên $n=2+8+15+10+5=40$.
\begin{center}
\begin{tabular}{|c|c|c|c|c|c|}
\hline
Chiều cao (cm) & $[10;20)$ & $[20;30)$ & $[30;40)$ & $[40;50)$ & $[50;60)$\\
\hline
Giá trị đại diện&$15$&$25$&$35$&$45$&$55$\\
\hline
Tần số & $2$ & $8$ & $15$ & $10$ & $5$\\
\hline
\end{tabular}
\end{center}
Độ tuổi trung bình của học viên trung tâm bằng
$$\overline{x}=\dfrac{15\cdot 2+25\cdot 8+35\cdot 15+45\cdot 10+55\cdot 5}{40}\approx 37$$
}
\end{bt}

\begin{ex}%[1D5H2-2]
% nguon: Archive/2024/Update-SP-BG10-11-3/SP-3/BG11-3/data_1C5/ChuongV-Bai1-Dang3-SoTrungVi-Diamon.tex #1523

Khi điều tra về chiều cao của $40$ học sinh lớp $11$A, bộ phận y tế của nhà trường thu được mẫu số liệu ghép nhóm như sau
\begin{center}
\begin{tabular}{|c|c|c|c|c|c|}
\hline
Chiều cao (cm)&$[140;146)$  &$[146;152)$  &$[152;158)$  &$[158;164)$  &$[164;170)$  \\
\hline
Số học sinh&$3$  &$12$  &$11$  &$8$  &$6$  \\
\hline
\end{tabular}
\end{center}
Xét tính đúng - sai của các khẳng định sau:
\choiceTF[t]
{Số học sinh có chiều cao trên $152$ cm là $24$}
{\True Tất cả học sinh của lớp đều có chiều cao không dưới $140$ cm}
{\True Nhóm chứa trung vị là nhóm thứ 3}
{\True Chiều cao trung bình của $40$ học sinh là $155{,}3$ cm}
\loigiai{
Ta có 
\begin{center}
\begin{tabular}{|c|c|c|c|c|c|}
\hline
Chiều cao (cm)&$[140;146)$  &$[146;152)$  &$[152;158)$  &$[158;164)$  &$[164;170)$  \\
\hline
Giá trị đại diện& $143$&$149$&$155$&$161$&$167$\\
\hline
Tần số&$3$  &$12$  &$11$  &$8$  &$6$  \\
\hline
\end{tabular}
\end{center}
\begin{itemchoice}
\itemch Só học sinh có chiều cao trên $152$ cm là $25$.
\itemch Tất cả học sinh của lớp đều có chiều cao không dưới $140$ cm.
\itemch Gọi $x_1; x_2; \ldots; x_{40}$ là chiều cao của $40$ học sinh nam sắp xếp theo thứ tự không giảm.\\
Trung vị của dãy số liệu là $\dfrac{x_{20} + x_{21}}{2}$. Mà $x_{20}$ và $x_{21}$ đều thuộc nhóm $[152;158)$ nên nhóm chứa trung vị là nhóm thứ $3$.
\itemch Chiều cao trung bình của học sinh nam trong lớp là
$$\overline{x} = \dfrac{143 \cdot 3 + 149 \cdot 12 + 155 \cdot 11 + 161 \cdot 8 + 167\cdot 3}{40} = 155{,}3\, \text{(cm)}.$$
\end{itemchoice}
}
\end{ex}

\begin{bt}%[1D5H2-2]
% nguon: Archive/2024/Update-SP-BG10-11-3/SP-3/BG11-3/data_1C5/ChuongV-Bai1-Dang3-SoTrungVi-Diamon.tex #1525
%[Duan11-1-GK1-NH23-24--TeamTeXHoa--HuynhThanhChi]
Khảo sát thời gian tham gia câu lạc bộ trong một tuần của  bạn học sinh, thu được kết quả như sau
\begin{center}
\begin{tabular}{|l|c|c|c|c|c|}
\hline
Thời gian (phút)& $[30;60)$ & $[60;90)$ & $[90;120)$ & $[120;150)$ & $[150;180)$ \\
\hline
Số học sinh& $36$ & $28$ & $18$ & $14$ & $4$\\
\hline
\end{tabular}
\end{center}
Số trung vị của mẫu số liệu trên là
\shortans[]{$75$}
\loigiai{
Cỡ mẫu $n=36+28+18+14+4=100$ và $\dfrac{n}{2}=50$.\\
Nhóm chứa trung vị là $[60;90)$. Suy ra $u_m=60$ và $u_{m+1}=90$.\\
Tần số của nóm chứa trung vị là $n_m=28$ và tần số tích lũy $C=n_1=36$.\\
Vậy trung vị của mẫu số liệu ghép nhóm là $M_e=60+\dfrac{50-36}{28}\cdot(90-60)=75$.
}
\end{bt}

\begin{ex}%[1D5N1-2]
% nguon: Archive/2024/Update-SP-BG10-11-3/SP-3/BG11-3/data_1C5/ChuongV-Bai1-Dang1-MauSoLieuGhepNhom-Diamon.tex #1531 
% KIEM_TRA: cần xem lại đáp số/dạng
%[Dự án EX-TF-2024, Mui Doan]
Trong một cuộc khảo sát đã tiến hành xác định tuổi (theo năm) của $120$ chiếc ô tô. Kết quả điều tra được cho trong bảng sau
\begin{center}
\begin{tabular}{|c|c|c|c|c|c|}
\hline Số tuổi (theo năm)&$[0;4)$&$[4;8)$&$[8;12)$&$[12;16)$&$[20;24)$\\
\hline
Số ô tô&$23$&$25$&$37$&$26$&$19$\\
\hline
\end{tabular}
\end{center}
Các khẳng định sau đúng hay sai?
\choiceTF[t]
{\True Số ô tô có độ tuổi dưới $12$ là $85$}
{Số ô tô có độ tuổi dưới $8$ là $23$}
{\True Số ô tô có độ tuổi dưới $16$ là $111$}
{\True Số ô tô có độ tuổi dưới $24$ là $130$}
\loigiai{ 
\begin{itemchoice}
\itemch  Số ô tô có độ tuổi dưới $12$ là $23+25+37=85$ ô tô.
\itemch  Số ô tô có độ tuổi dưới $8$ là $23+25=48$ ô tô.
\itemch  Số ô tô có độ tuổi dưới $16$ là $23+25+37+26=111$ ô tô.
\itemch  Số ô tô có độ tuổi dưới $24$ là $23+25+37+26+19=130$ ô tô.
\end{itemchoice}

}
\end{ex}

\begin{ex}%[1D5N1-2]
% nguon: Archive/2024/Update-SP-BG10-11-3/SP-3/BG11-3/data_1C5/ChuongV-Bai1-Dang1-MauSoLieuGhepNhom-Diamon.tex #1532
%[Dự án EX-TF-2024, Mui Doan]
Điều tra về chiều cao của $100$ học sinh lớp $11$ của một trường THPT, ta được kết quả
\begin{center}
\begin{tabular}{|c|c|c|c|c|c|c|c|}
\hline \begin{tabular}{c} 
Chiều cao \\
(cm)
\end{tabular} & {$[150 ; 152)$} & {$[152 ; 154)$} & {$[154 ; 156)$} & {$[156 ; 158)$} & {$[158 ; 160)$} & {$[160 ; 162)$} & {$[162 ; 168)$} \\
\hline Số học sinh & $5$ & $18$ & $40$ & $25$ & $8$ & $3$ & $1$ \\
\hline
\end{tabular}
\end{center}
Các khẳng định sau đúng hay sai?
\choiceTF[t]
{\True Mẫu số liệu trên có $7$ nhóm}
{Số học sinh có chiều cao dưới $154$ cm là $18$}
{\True Số học sinh có chiều cao trên $160$ cm là $4$}
{\True Số học sinh có chiều cao trong khoảng từ $156$ cm đến dưới $160$ cm là $33$}
\loigiai{
\begin{itemchoice}
\itemch  Mẫu số liệu trên có $7$ nhóm.
\itemch  Số học sinh có chiều cao dưới $154$ cm là $5+18=23$.
\itemch   Số học sinh có chiều cao trên $160$ cm là $3+1=4$.
\itemch  Số học sinh có chiều cao trong khoảng từ $156$ cm đến dưới $160$ cm là $25+8=33$.
\end{itemchoice}
}
\end{ex}

\begin{ex}%[1D5N1-2]
% nguon: Archive/2024/Update-SP-BG10-11-3/SP-3/BG11-3/data_1C5/ChuongV-Bai1-Dang1-MauSoLieuGhepNhom-Diamon.tex #1533
%[Dự án EX-TF-2024, Mui Doan]
Cho bảng số liệu khảo sát về tuổi thọ (đơn vị: nghìn giờ) của một loại bóng đèn:
\begin{longtable}{|c|c|c|c|c|c|}
\hline Tuổi thọ &$[3 ; 5)$&$[5 ; 7)$&$[7 ; 9)$&$[9 ; 11)$&$[11 ; 13)$\\
\hline Số bóng đèn & $4$ & $20$ & $26$ & $42$ & $8$ \\
\hline
\end{longtable}
Các khẳng định sau đúng hay sai? 
\choiceTF[t]
{Số bóng đèn có tuổi thọ dưới $7$ nghìn giờ là $4$}
{\True Số bóng đèn có tuổi thọ dưới $9$ nghìn giờ là $50$}
{Số bóng đèn có tuổi thọ dưới $13$ nghìn giờ là $92$}
{\True Số bóng đèn có tuổi thọ từ $9$ nghìn giờ trở lên là $50$}
\loigiai{
\begin{itemchoice}
\itemch  Số bóng đèn có tuổi thọ dưới $7$ nghìn giờ là $4+20=24$.
\itemch  Số bóng đèn có tuổi thọ dưới $9$ nghìn giờ là $4+20+26=50$.
\itemch  Số bóng đèn có tuổi thọ dưới $13$ nghìn giờ là $4+20+26+42+8=100$.
\itemch  Số bóng đèn có tuổi thọ từ $9$ nghìn giờ trở lên là $42+8=50$.
\end{itemchoice}
}
\end{ex}

\begin{ex}%[1D5N1-2]
% nguon: Archive/2024/Update-SP-BG10-11-3/SP-3/BG11-3/data_1C5/ChuongV-Bai1-Dang1-MauSoLieuGhepNhom-Diamon.tex #1534
%[Dự án EX-TF-2024, Mui Doan]
Số lượng khách hàng nữ mua hàng thời trang trong một ngày của một cửa hàng được thống kê trong bảng tần số ghép nhóm sau

\begin{center}
\begin{tabular}{|c|c|c|c|c|c|}
\hline Khoảng tuổi & $[20 ; 30)$ & $[30 ; 40)$ & $[40 ; 50)$ & $[50 ; 60)$ & $[60 ; 70)$ \\
\hline Số khách hàng nữ & $3$ & $9$ & $6$ & $4$ & $2$ \\
\hline
\end{tabular}
\end{center}
Các khẳng định sau đúng hay sai? 
\choiceTF[t]
{Số khách hàng nữ có độ tuổi dưới $40$ là $ 3$}
{\True Số khách hàng nữ có độ tuổi trên  $50$ là $ 6$}
{\True Số khách hàng nữ có độ tuổi từ $30$ đến dưới $40$ là $9$}
{\True Cỡ mẫu của bảng tần số ghép nhóm là $ 24$}
\loigiai{
\begin{itemchoice}
\itemch  Số khách hàng nữ có độ tuổi dưới $40$ là $ 3+9=12$.
\itemch  Số khách hàng nữ có độ tuổi trên  $50$ là $ 4+2=6$.
\itemch  Số khách hàng nữ có độ tuổi từ $30$ đến dưới $40$ là $9$.
\itemch  Cỡ mẫu của bảng tần số ghép nhóm là $ 3+9+6+4+2=24$.
\end{itemchoice}
}
\end{ex}

\begin{ex}%[1D5H1-2]
% nguon: Archive/2024/Update-SP-BG10-11-3/SP-3/BG11-3/data_1C5/ChuongV-Bai1-Dang1-MauSoLieuGhepNhom-Diamon.tex #1535
%[Dự án True-False-2024]%[Phan Văn Thành]
Cho mẫu số liệu về thời gian (phút) đi từ nhà đến nơi làm việc  của các nhân viên một công ty như sau
\begin{center}
\begin{tabular}{|c|c|c|c|c|c|c|c|}
\hline
Thời gian (phút)&$[15;20)$  &$[20;25)$  &$[25;30)$  &$[30;35)$  &$[35;40)$&$[40;45)$& $[45;50)$  \\
\hline
Số nhân viên&$6$  &$14$  &$25$  &$37$  &$21$& $13$& $9$  \\
\hline
\end{tabular}
\end{center}
Xét tính đúng - sai của các khẳng định sau:
\choiceTF[t]
{\True Mẫu số liệu trên được ghép thành $7$ nhóm}
{Độ dài mỗi nhóm của mẫu số liệu bằng $15$}
{\True Số nhân viên đi từ nhà đến nơi làm việc từ $25$ phút đến dưới $30$ phút là $25$}
{Nhóm có số nhân viên nhiều nhất là nhóm $[45;50)$}
\loigiai{
\begin{itemchoice}
\itemch Mẫu số liệu được ghép thành $7$ nhóm.
\itemch Độ dài mỗi nhóm của mẫu số liệu là bằng $5$.
\itemch Số nhân viên từ nhà đến nơi làm việc từ $25$ phút đến dưới $30$ phút là $25$.
\itemch Nhóm có số nhân viên nhiều nhất là nhóm $[30;35)$.
\end{itemchoice}
}
\end{ex}

\begin{bt}%[1D5N1-2]
% nguon: Archive/2024/Update-SP-BG10-11-3/SP-3/BG11-3/data_1C5/ChuongV-Bai1-Dang1-MauSoLieuGhepNhom-Diamon.tex #1537

Khảo sát tổng thời gian truy cập internet mỗi tối (đơn vị: phút) của một số học sinh thu được mẫu số liệu ghép nhóm sau
\begin{center}
\begin{tabular}{|c|c|c|c|c|c|c|}
\hline
Thời gian (phút)& $\left[0;30\right)$ & $\left[30;60\right)$  & $\left[60;90\right)$&$\left[90;120\right)$ &$\left[120;150\right)$  &$\left[150;180\right)$  \\
\hline
Số học sinh& $11$ & $25$ &$36$  & $15$ & $8$ &$5$ \\
\hline
\end{tabular}
\end{center}
Tổng các giá trị đại diện của mẫu số liệu ghép nhóm trên bằng
\shortans[]{$540$}
\loigiai{
Ta có
\begin{center}
\begin{tabular}{|c|c|c|c|c|c|c|}
\hline
Thời gian (phút)& $\left[0;30\right)$ & $\left[30;60\right)$  & $\left[60;90\right)$&$\left[90;120\right)$ &$\left[120;150\right)$  &$\left[150;180\right)$  \\
\hline
Giá trị đại diện& $15$ & $45$ &$75$  & $105$ & $135$ &$165$ \\
\hline
\end{tabular}
\end{center}
Do đó, tổng các giá trị đại diện bằng $15+45+75+105+135+165=540$.
}
\end{bt}

\begin{bt}%[1D5N1-2]
% nguon: Archive/2024/Update-SP-BG10-11-3/SP-3/BG11-3/data_1C5/ChuongV-Bai1-Dang1-MauSoLieuGhepNhom-Diamon.tex #1538
 
Khảo sát thời gian tự học trong ngày (đơn vị: giờ) của một số học sinh khối $11$ thu được mẫu số liệu ghép nhóm sau
\begin{center}
\begin{tabular}{|c|c|c|c|c|}
\hline
Thời gian (giờ)& $\left[0;2\right)$ & $\left[2;4\right)$  & $\left[4;6\right)$&$\left[6;8\right)$ \\
\hline
Số học sinh& $6$ & $21$ &$28$  & $5$ \\
\hline
\end{tabular}
\end{center}
Tổng tần số tích lũy của nhóm $[ 4;6 )$ là
\shortans[]{$55$}
\loigiai{
Tổng tần số tích lũy của nhóm $[ 4;6 )$ là $6+21+28=55$.
} 
\end{bt}

\begin{ex}%[1D5H2-3]
% nguon: Archive/2024/Update-SP-BG10-11-3/SP-3/BG11-3/data_1C5/DeOT-ChuongV-DeOnTap-So1-VuNguyenHoangAnh.tex #1579

	Mẫu số liệu dưới đây ghi lại độ dài quãng đường di chuyển trong một tuần (đơn vị: kilômét) của $40$ chiếc ô tô
	\begin{center}
		\begin{tabular}{|c|c|c|}
			\hline
			Nhóm & Tần số & Tần số tích lũy \\
			\hline
			$[100;120)$ & $4$ & $4$ \\
			\hline
			$[120;140)$ & $10$ & $14$ \\
			\hline
			$[140;160)$ & $19$ & $33$ \\
			\hline
			[160;180) & $5$ & $38$ \\
			\hline
			$[180;200)$ & $2$ & $40$ \\
			\hline
			\multicolumn{3}{|c|}{$n=40$} \\
			\hline
		\end{tabular}
	\end{center}
	\choiceTF[t]
	{\True Trung vị của mẫu số liệu thuộc nhóm $[140;160)$}
	{\True Số trung bình của mẫu số liệu là $145{,}5$}
	{\True Mốt của mẫu số liệu thuộc nhóm $[140;160)$}
	{Hiệu của hai tứ phân vị thứ ba và thứ nhất không quá $20$}
	\loigiai{
		\begin{itemchoice}
			\itemch Ta thấy $x_{20}, x_{21}\in [140;160)$ nên $M_e$ thuộc nhóm này.
			\itemch Số trung bình của mẫu số liệu là 
			\[ \overline{x}=\frac{4\cdot 110+10\cdot 130+19\cdot 150+5\cdot 170+2\cdot 190}{40}=145{,}5. \]
			\itemch Ta thấy nhóm $3$ ứng với nửa khoảng $[140;160)$ là nhóm có tần số lớn nhất là $19$, hai nhóm trước và sau lần lượt có tần số là $10$ và $5$. Áp dụng công thức, ta có mốt của mẫu số liệu là
			$$ M_o=140+\left(\dfrac{19-10}{2\cdot 19-10-5}\right) \cdot 20 \approx 148.$$
			\itemch \begin{itemize} 
				\item Ta có $\dfrac{n}{4}=10$. Nhóm $2$ là nhóm đầu tiên có tần số tích lũy lớn hơn hoặc bằng $10$. Áp dụng công thức, ta có tứ phân vị thứ nhất là
				$$Q_1=120+\frac{10-4}{10}\cdot 20 =132.$$
				\item Ta có $\dfrac{3n}{4}=30$. Nhóm $3$ là nhóm đầu tiên có tần số tích lũy lớn hơn hoặc bằng $30$. Áp dụng công thức, ta có tứ phân vị thứ ba là
				$$Q_3=140+\left(\dfrac{30-14}{19}\right)\cdot 20 \approx 157.$$	
			\end{itemize}
			Do đó $Q_3-Q_1\approx 25$. 
		\end{itemchoice} 
	}
\end{ex}

\begin{bt}%[1D5N1-3]
% nguon: Archive/2024/Update-SP-BG10-11-3/SP-3/BG11-3/data_1C5/DeOT-ChuongV-DeOnTap-So1-VuNguyenHoangAnh.tex #1582

		Chiều cao của $20$ cây giống  được cho bởi bảng sau
		\begin{center}
			\begin{tabular}{|c|c|}
				\hline  
				{\bf Lớp} (cm)&{\bf Tần số}\\
				\hline 
				$[40;44]$ & $2$ \\
				$[45;49]$ & $5$ \\
				$[50;54]$ & $3$ \\
				$[55;59]$ & $4$ \\
				$[60;64]$ & $3$ \\
				$[65;69]$ & $3$ \\
				\hline
				& $N=20$ \\
				\hline
			\end{tabular} 
		\end{center}
		Số trung bình $\bar{x}$ (làm tròn đến chữ số thứ hai sau dấu phẩy) của bảng nói trên là bao nhiêu?
		\shortans{54,5}
		\loigiai{
			Bảng tần số ghép lớp của bảng nói trên là
			\begin{center}
				\begin{tabular}{|c|c|c|}
					\hline  
					{\bf Lớp} (cm)& Giá trị đại diện &{\bf Tần số}\\
					\hline 
					$[40;44]$ &$42$& $2$ \\
					$[45;49]$ &$47$& $5$ \\
					$[50;54]$ &$52$& $3$ \\
					$[55;59]$ & $57$&$4$ \\
					$[60;64]$ & $62$&$3$ \\
					$[65;69]$ &$67$ &$3$ \\
					\hline
					&& $N=20$ \\
					\hline
				\end{tabular} 
			\end{center}
			Áp dụng công thức tính số trung bình cho bảng tần số ghép lớp ta có
			\begin{center}
				$\bar{x} \approx \dfrac{42\cdot 2+47\cdot 5+52\cdot 3+57\cdot 4+62\cdot 3+67\cdot3}{20} = 54{,}5$ (cm).
			\end{center}
		}
\end{bt}

\begin{bt}%[1D5N1-4]
% nguon: Archive/2024/Update-SP-BG10-11-3/SP-3/BG11-3/data_1C5/DeOT-ChuongV-DeOnTap-So1-VuNguyenHoangAnh.tex #1583

	Thống kê điểm trung bình môn Toán của một số học sinh lớp $11$ được cho ở bảng sau:
	\begin{center}
		\begin{tabular}{|c|c|c|c|c|c|c|c|}
			\hline Khoảng điểm &{$[6{,}5; 7)$} &{$[7; 7{,}5)$} &{$[7{,}5; 8)$} &{$[8; 8{,}5)$} &{$[8{,}5;9)$} &{$[9; 9{,}5)$} &{$[9{,}5; 10)$} \\
			\hline Tần số & $8$ & $10$ & $16$ & $24$ & $13$ & $7$ & $4$ \\
			\hline
		\end{tabular}
	\end{center}
	Mốt của mẫu số liệu trên là bao nhiêu (làm tròn đến hai chữ số sau dấu phẩy)?
	\shortans{8,21}
	\loigiai{
		Nhóm chứa mốt của mẫu số liệu trên là nhóm $[8; 8{,}5)$.\\
		Do đó $u_m=8$; $n_{m-1}=16$; $n_m=24$; $n_{m+1}=13$; $u_{m+1}-u_m=8{,}5-8=0{,}5$.\\
		Mốt của mẫu số liệu ghép nhóm là
		\[M_o=8+\dfrac{24-16}{(24-16)+(24-13)} \cdot 0{,}5\approx 8{,}21. \] 
	}
\end{bt}

\begin{bt}%[1D5V2-3]
% nguon: Archive/2024/Update-SP-BG10-11-3/SP-3/BG11-3/data_1C5/DeOT-ChuongV-DeOnTap-So1-VuNguyenHoangAnh.tex #1584

	Một bảng xếp hạng đã tính điểm chuẩn hoá cho chỉ số nghiên cứu của một số trường đại học ở Việt Nam và thu được kết quả như sau
	\begin{center}
		\begin{tabular}{|c|c|c|c|c|c|c|}
			\hline 
			Điểm  & Dưới $20$ &$[20;30)$&$[30;40)$&$[40;60)$&  $[60;80)$&$[80;100)$\\
			\hline 
			Số lượng&$4$&$19$&$6$&$2$&$3$&$1$\\
			\hline 
		\end{tabular}
	\end{center}	
	Xác định điểm ngưỡng để đưa ra danh sách $25\%$ trường đại học có chỉ số nghiên cứu tốt nhất Việt Nam (lam tròn đến một chữ số thập phân).
	\shortans{35,4}
	\loigiai{
		Điểm ngưỡng để đưa ra danh sách $25\%$ trường đại học có chỉ số nghiên cứu tốt nhất Việt Nam là tứ phân vị thứ ba $Q_3$.\\
		Mẫu số liệu có $35$ trường và do $\dfrac{3\cdot 35}{4}=26{,}25$ nên tứ phân vị thứ ba thuộc nhóm thứ ba. Do đó
		$$Q_3=30+\dfrac{\dfrac{3\cdot 35}{4}-4-19}{6}\cdot 10\approx 35{,}4.$$
	} 
\end{bt}

\begin{bt}%[1D5V1-4]
% nguon: Archive/SP_KhoLatex/ID6_old/1-CK1-OT-De-34.tex #2708

Số cuộc gọi điện thoại một người thực hiện mỗi ngày trong $30$ ngày được lựa chọn ngẫu nhiên và được thống kê trong bảng sau:
\begin{center}
\begin{tabular}{|c|c|c|c|c|c|}
	\hline Số cuộc gọi & {$[2{,}5 ; 5{,}5)$} & {$[5{,}5 ; 8{,}5)$} & {$[8{,}5 ; 11{,}5)$} & {$[11{,}5 ; 14{,}5)$} & {$[14{,}5 ; 17{,}5)$} \\
	\hline Số ngày & 5 & 13 & 7 & 3 & 2 \\
	\hline
\end{tabular}
\end{center}
Hãy xác định mốt của mẫu số liệu ghép nhóm trên.
\loigiai{
Ta có nhóm chứa mốt của mẫu số liệu trên là nhóm $\left[5{,}5;8{,}5\right)$.\\
Do đó $u_m=5,5$, $n_{m-1}=5$, $n_m=13$, $n_{m+1}=7$, $u_{m+1}-u_m=8{,}5-5{,}5=3$.\\
Mốt của mẫu số liệu ghép nhóm là: $M_0=5{,}5+\dfrac{13-5}{(13-5)+(13-5)}\cdot 3=\dfrac{101}{14}\approx 7{,}2$.}
\end{bt}

\begin{bt}%[1D5H1-3]
% nguon: Archive/SP_KhoLatex/ID6_old/1-CK1-OT-De-32.tex #2806
%[Dự án EX-CK1 10 và 11 mới - Thành Đức Trung]
	Thống kê điểm trung bình môn Toán của một số học sinh lớp $11$ được cho ở bảng sau:
	\begin{center}
		\begin{tabular}{|c|c|c|c|c|c|c|c|}
			\hline
			Khoảng điểm &$\left[6{,}5;7\right)$&$\left[7;7{,}5\right)$&$\left[7{,}5;8\right)$&$\left[8;8{,}5\right)$&$\left[8{,}5;9\right)$&$\left[9;9{,}5\right)$&$\left[9{,}5;10\right)$\\
			\hline
			Tần số&$8$&$10$&$16$&$24$&$13$&$7$&$4$\\
			\hline
		\end{tabular}
	\end{center}
	Số trung bình của mẫu số liệu trên bằng bao nhiêu?
	\loigiai{
		Ta có bảng sau
		\begin{center}
			\begin{tabular}{|c|c|c|c|c|c|c|c|}
				\hline
				Khoảng điểm &$\left[6{,}5;7\right)$&$\left[7;7{,}5\right)$&$\left[7{,}5;8\right)$&$\left[8;8{,}5\right)$&$\left[8{,}5;9\right)$&$\left[9;9{,}5\right)$&$\left[9{,}5;10\right)$\\
				\hline
				Giá trị đại diện &$6{,}75$&$7{,}25$&$7{,}25$&$8{,}25$&$8{,}75$&$9{,}25$&$9{,}75$\\
				\hline
				Tần số&$8$&$10$&$16$&$24$&$13$&$7$&$4$\\
				\hline
			\end{tabular}
		\end{center}
		Khi đó $\overline{x}=\dfrac{8\cdot 6{,}75+10\cdot 7{,}25+16\cdot 7{,}75+24\cdot 8{,}25+13\cdot 8{,}75+7\cdot 9{,}25+4\cdot 9{,}75}{82}=\dfrac{333}{41}\approx 8{,}1$.}
\end{bt}

\begin{bt}%[1D5V2-2]
% nguon: Archive/SP_KhoLatex/ID6_old/1-CK1-OT-De-29.tex #3652
%Câu 40
	(0.5 điểm) Quãng đường (km) các cầu thủ (không tính thủ môn) chạy trong một trận bóng đá tại giải ngoại hạng Anh được cho trong bảng thống kê sau\\
	\begin{center}
		\begin{tabular}{|l|c|c|c|c|c|}
			\hline Quãng đường & $[2 ; 4)$ & $[4 ; 6)$ & $[6 ; 8)$ & $[8 ; 10)$ & $[10 ; 12)$ \\
			\hline  Số cầu thủ  & $2$ & $5$ & $6$ & $9$ & $3$ \\
			\hline
		\end{tabular}
	\end{center}
	Tìm trung vị của mẫu số liệu và giải thích ý nghĩa của giá trị thu được.
\loigiai{
	Cỡ mẫu $n=2+5+6+9+3=25$.\\
	Gọi $x_1$, $x_2$, $\ldots$, $x_{25}$ là quãng đường chạy của $25$ cầu thủ và giả sử dãy này sắp xếp theo thứ tự không giảm.\\
	Khi đó trung vị là $x_{13}$, nên nhóm chứa trung vị là $\left[6;8\right)$. \\
	Trung vị là\\
	$M_e=6+\dfrac{\dfrac{25}{2}-\left(2+5\right)}{6}\cdot \left(8-6\right)\approx 7{,}83$\\
	Có $50\%$ số cầu thủ chạy nhiều hơn $7{,}83$ km và có $50\%$ số cầu thủ chạy ít hơn $7{,}83$ km.\\
}
\end{bt}

\begin{bt}%[1D5H1-4]
% nguon: Archive/SP_KhoLatex/SP_Đề Theo Bài-Khối 11/1-B8.tex #4152

	Khảo sát số lần sử dụng Facebook của một người thực hiện mỗi ngày trong $30$ ngày được lựa chọn ngẫu nhiên được thống kê trong bảng sau:
\begin{center}
\begin{tabular}{|c|c|c|c|c|c|}
\hline \begin{tabular}{c} 
Số lần sử \\
dụng Facebook
\end{tabular} & {$[3 ; 5]$} & {$[6 ; 8]$} & {$[9 ; 11]$} & {$[12 ; 14]$} & {$[15 ; 17]$} \\
\hline Số ngày & $2$ & $5$ & $11$ & $8$ & $4$ \\
\hline
\end{tabular}
\end{center}
Tìm mốt của mẫu số liệu ghép nhóm trên?
\shortans{$10{,}5$}
\loigiai{
Do số lần sử dụng Facebook là số nguyên nên ta chỉnh lại như sau:
\begin{center}
\begin{tabular}{|c|c|c|c|c|c|}
\hline \begin{tabular}{c} 
Số lần sử \\
dụng Facebook
\end{tabular} & {$[2,5 ; 5,5)$} & {$[5,5 ; 8,5)$} & {$[8,5 ; 11,5)$} & {$[11,5 ; 14,5)$} & {$[14,5 ; 17,5)$} \\
\hline Số ngày & $2$ & $5$ & $11$ & $8$ & $4$ \\
\hline
\end{tabular}
\end{center}
Nhóm chứa mốt của mẫu số liệu trên là $[8{,}5 ; 11{,}5)$.\\
Do đó: $u_m=8{,}5 ;$ $ n_{m-1}=5 ;$ $ n_{m+1}=8 ;$ $ u_{m+1}-u_m=11{,}5-8{,}5=3$.\\
Mốt của mẫu số liệu ghép nhóm là $M_o=8{,}5+\dfrac{11-5}{(11-5)+(11-8)} \cdot 3=10{,}5$.}
\end{bt}

\begin{bt}%[1D5H1-4]
% nguon: Archive/SP_KhoLatex/SP_Đề Theo Bài-Khối 11/1-B8.tex #4155

	Công ty bất động sản Đất Vàng thực hiện cuộc khảo sát khách hàng xem họ có nhu cầu mua nhà ở mức giá nào để tiến hành dự án xây nhà ở Thăng Long group sắp tới. Kết quả khảo sát $500$ khách hàng được ghi lại ở bảng sau:
\begin{center}
\begin{tabular}{|c|l|l|l|l|l|}
\hline \begin{tabular}{c} 
Mức giá (triệu \\
đồng)
\end{tabular} & {$[10 ; 14)$} & {$[14 ; 18)$} & {$[18 ; 22)$} & {$[22 ; 26)$} & {$[26 ; 30)$} \\
\hline Số khách hàng & 75 & 105 & 179 & 96 & 45 \\
\hline
\end{tabular}
\end{center}
Công ty bất động sản Đất Vàng nên xây nhà ở mức giá nào để nhiều người có nhu cầu xây nhà?
\shortans{$19{,}9$}
\loigiai{
Nhóm chứa mốt của mẫu số liệu trên là $[18;22)$.\\
Do đó: $u_m=18 ;$ $ n_{m-1}=105 ;$ $ n_{m+1}=96 ;$ $ u_{m+1}-u_{m}=22-18=4$.\\
Mốt của mẫu số liệu ghép nhóm là $M_o=18+\dfrac{179-105}{(179-105)+(179-96)} \cdot 4 \approx 19{,}9$.\\
Dựa vào kết quả trên ta có thể dự đoán rằng nếu công ty bất động sản Đất Vàng xây nhà ở mức giá $19{,}9$ triệu đồng $/ \text{m}^2$ thì sẽ có nhiều người mua nhất.
}
\end{bt}

\begin{bt}%[1D5H1-4]
% nguon: Archive/SP_KhoLatex/SP_Đề Theo Bài-Khối 11/1-B8.tex #4157

	Nghiên cứu thời gian chạy $01$ vòng sân trường $300 \mathrm{~m}$ của $41$ học sinh lớp $11 \mathrm{~A}$ trường THPT được giáo viên bộ môn Thể dục ghi lại, có kết quả sau:
\begin{center}
\begin{tabular}{|c|c|c|c|c|c|}
\hline Thời gian & {$[40 ; 45)$} & {$[45 ; 50)$} & {$[50 ; 55)$} & {$[55 ; 60)$} & {$[60 ; 65)$} \\
\hline Số học sinh & $5$ & $8$ & $13$ & $9$ & $6$ \\
\hline
\end{tabular}
\end{center}
Hãy tìm mốt của mẫu số liệu ghép nhóm trên.
\shortans{$52{,}5$}
\loigiai{
Nhóm chứa mốt của mẫu số liệu trên là $[50 ; 55 )$.\\
Do đó: $u_m=50 ;$ $ n_{m-1}=8 ;$ $ n_{m+1}=9 ;$ $ u_{m+1}-u_m=55-50=5$.\\
Mốt của mẫu số liệu ghép nhóm là $$M_o=50+\dfrac{13-9}{(13-9)+(13-9)} \cdot 5=52{,}5.$$
}
\end{bt}

\begin{bt}%[1D5V2-3]
% nguon: Archive/2025-Sangtac/Khoi10/HK2/0-HK2-2025-SangTac/Data/0-CTST/CTST_2-OTHK1_De04.tex #4439

	Biểu đồ dưới đây biểu diễn số lượt khách hàng đặt bàn qua hình thức trực tuyến mỗi ngày trong quý III năm $2022$ của một nhà hàng. Cột thứ nhất biểu diễn số ngày có từ $1$ đến dưới $6$ lượt đặt bàn; cột thứ hai biểu diễn số ngày có từ $6$ đến dưới $11$ lượt đặt bàn;\ldots
	\begin{center}
		\begin{tikzpicture}
			\begin{axis}[
				width=12cm, % Độ rộng biểu đồ
				height=8cm, % Chiều cao biểu đồ
				ybar,
				bar width=0.8, % Độ rộng của cột
				enlargelimits=0.15,
				ymin=0, ymax=35, % Thay đổi ymax để phù hợp với dữ liệu
				xmin=-2, xmax=30, % Thay đổi xmax để phù hợp với trục x
				xtick={1,6,11,16,21,26}, % Các vị trí tick trên trục x (sửa cho phù hợp với các khoảng)
				ytick={0,5,10,15,20,25,30,35}, % Các vị trí tick trên trục y
				xlabel={Số lượt đặt bản}, % Nhãn trục x
				ylabel={Số ngày}, % Nhãn trục y
				title={Số lượt khách hàng đặt bàn trực tuyến}, % Tiêu đề
				ymajorgrids=true, % Đường lưới chính trên trục y
				grid style=dashed, % Kiểu đường lưới
				axis x line=bottom, % Định dạng trục x
				axis y line=left, % Định dạng trục y
				bar width=5pt % Độ rộng của cột
				]
				
				% Dữ liệu biểu đồ
				\addplot[
				ybar,
				bar width=46pt, % Độ rộng của cột (điều chỉnh độ rộng để phù hợp với các khoảng)
				fill=gray
				] 
				coordinates {
					(3.5,14) (8.5,30) (13.5,25) (18.5,18) (23.5,5)
				};
				
				% Nhãn các giá trị trên cột
				\node at (axis cs:1.5,14)[above right,font=\small] {14};
				\node at (axis cs:6.5,30)[above right,font=\small] {30};
				\node at (axis cs:11.5,25)[above right,font=\small] {25};
				\node at (axis cs:16.5,18)[above right,font=\small] {18};
				\node at (axis cs:21.5,5)[above right,font=\small] {5};
				
			\end{axis}
		\end{tikzpicture}
	\end{center}
	Hãy tìm khoảng tứ phân vị của mẫu số liệu ghép nhóm cho bởi biểu đồ trên.
	\shortans{$8{,}5$}
	\loigiai{
		Cỡ mẫu $n=92$.\\
		Gọi $x_1; x_2; \ldots; x_{92}$ là mẫu số liệu gốc về số lượt khách hàng đặt bàn qua hình thức trực tuyến mỗi ngày trong quý III năm 2022 của nhà hàng được xếp theo thứ tự không giảm.\\
		Ta có:
		\allowdisplaybreaks
		\begin{align*}
			&x_1, \ldots, x_{14} \in [1; 6); \\
			&x_{15}, \ldots, x_{44} \in [6; 11); \\
			&x_{45}, \ldots, x_{69} \in [11; 16); \\
			&x_{70}, \ldots, x_{87} \in [16; 21); \\
			&x_{88}, \ldots, x_{92} \in [21; 26).
		\end{align*}
		Tứ phân vị thứ nhất của mẫu số liệu gốc là $\dfrac{1}{2}\left(x_{23} + x_{24}\right) \in [6; 11)$.\\
		Do đó, tứ phân vị thứ nhất của mẫu số liệu ghép nhóm là:
		$$
		Q_1 = 6 + \dfrac{\dfrac{92}{4} - 14}{30} \cdot (11 - 6) = 7{,}5.
		$$
		Tứ phân vị thứ ba của mẫu số liệu gốc là $\dfrac{1}{2}\left(x_{69} + x_{70}\right) \in [16; 21)$.\\
		Do đó, tứ phân vị thứ ba của mẫu số liệu ghép nhóm là
		$Q_3 = 16$.\\
		Vậy khoảng tứ phân vị của mẫu số liệu ghép nhóm là
		$$
		\Delta_Q = Q_3 - Q_1 = 16 - 7{,}5 = 8{,}5.
		$$
	}
\end{bt}

\begin{bt}%[1D5V1-2]
% nguon: Archive/2025-Sangtac/Khoi10/Data/1-KNTT/KNTT_1-OTHK1_De02.tex #4445

	Thống kê thời gian hoàn thành bài kiểm tra Toán 90 phút của các bạn trong lớp và thu được mẫu số liệu sau:
	\begin{center}
		$\begin{array}{lllllllllllll}
			&89	&77	&78	&74	&79	&79	&89	&85	&86 &87	&71	&76\\
			&80	&87	&88	&89	&87	&81	&88	&86	&87 &70	&85	&71\\
			&78	&76	&86	&72	&88	&76	&72	&84	&71 &77	&83	&78\\
		 \end{array}$
	\end{center}
	Hãy chuyển mẫu số liệu trên sang dạng ghép nhóm (chia thành $4$ nhóm). Tính tỉ lệ của nhóm có số học sinh nộp bài nhiều nhất (làm tròn đến chữ số thập phân thứ hai).
	\shortans[0]{$0{,}42$}
	\loigiai{
	Bảng tần số ghép nhóm
		\begin{center}
		\begin{tabular}{|>{\centering\arraybackslash}p{3.5cm}|>{\centering\arraybackslash}p{2cm}|>{\centering\arraybackslash}p{2cm}|>{\centering\arraybackslash}p{2cm}|>{\centering\arraybackslash}p{2cm}|>{\centering\arraybackslash}p{2cm}|}
			\hline
			\textbf{Thời gian (phút)}&$[70 ; 75)$&$[75 ; 80)$&$[80 ; 85)$&$[85 ; 90)$\\
			\hline
			\textbf{Số học sinh} &$7$ &$10$ &$4$ &$15$\\
			\hline
		\end{tabular}
	\end{center}
	Cỡ mẫu $n=36$.\\
	Tỉ lệ của nhóm có số học sinh nộp bài nhiều nhất là $\dfrac{15}{36} \approx 0{,}42$.
	}
\end{bt}

\begin{bt}%[1D5V1-4]
% nguon: Archive/2025-Sangtac/Khoi10/Data/1-KNTT/KNTT_1-OTHK1_De02.tex #4446

	Mức thưởng tết (triệu đồng) mà nhân viên của một công ty nhận được bằng mẫu số liệu ghép nhóm sau:
	\begin{center}
		\begin{tabular}{|c|c|c|c|c|c|}
			\hline 
			\textbf{Mức thưởng}	& $ \left[ 5 ; 10\right)  $ & $ \left[ 10 ; 15\right)  $ & $ \left[ 15 ; 20\right)  $ & $ \left[ 20 ; 25\right)  $ & $ \left[25 ; 30 \right)  $ \\ 
			\hline 
			\textbf{Số nhân viên}	& $ 12 $ & $ 25 $ & $ 40 $ & $ 20 $ & $ 3 $ \\ 
			\hline 
		\end{tabular} 
	\end{center}
	Số nhân viên nhận được mức thưởng tết khoảng bao nhiêu triệu đồng là nhiều nhất?
		\shortans[0]{$17{,}1$}
	\loigiai{
		Cỡ mẫu $n=12+25+40+20+3=100$.\\
		Số nhân viên có mức thưởng tết từ 15 đến dưới 20 triệu đồng là nhiều nhất nên nhóm chứa mốt là nhóm $\left[ 15 ; 20\right)$.\\
		Ta có, $a_j=15$; $m_j=40$; $m_{j-1}=25$; $m_{j+1}=20$; $h=5$.\\
		Do đó, mốt của mẫu số liệu ghép nhóm là $$M_o=15+\dfrac{(40-25)}{(40-25)+(40-20)} \cdot 5 \approx 17{,}1.$$
		Vậy số nhân viên nhận được mức thưởng tết khoảng $17{,}1$ triệu đồng là nhiều nhất.
		}
\end{bt}

\begin{ex}%[1D5H2-3]
% nguon: Archive/2025-Sangtac/Khoi12/Data/2-KNTT/KNTT_2-OTHK1_De01.tex #5037

	\immini{
		\textbf{\textit{Bảng 5}} biểu diễn mẫu số liệu ghép nhóm về số tiền (đơn vị: nghìn đồng) mà $60$ khách hàng mua sách ở một cửa hàng trong một ngày
	}{
		\begin{tabular}{|c|c|c|}
			\hline
			Nhóm & Tần số & Tần số tích lũy\\
			\hline
			$[40;50)$ &$5$&$5$\\
			\hline
			$[50;60)$ &$8$&$13$\\
			\hline
			$[60;70)$ &$25$&$38$\\
			\hline
			$[70;80)$ &$20$&$58$\\
			\hline
			$[80;90)$ &$2$&$60$\\
			\hline
			& $n=60$&\\
			\hline
		\end{tabular}\\
		\centering {\textit{Bảng 5}}
	}
	\choiceTF
	{Số trung bình của mẫu số liệu trên là $65$ (nghìn đồng)}
	{\True Trung vị của mẫu số liệu trên là $66{,}8$ (nghìn đồng)}
	{\True Tứ phân vị thứ nhất $Q_1$ của mẫu số liệu trên là $60{,}8$ (nghìn đồng)}
	{Mốt của mẫu số liệu trên là $65$ (nghìn đồng)}
	\loigiai{
		Ta có
		\begin{center}
			\begin{tabular}{|c|c|c|c|c|c|}
				\hline
				Nhóm & $[40;50)$& $[50;60)$ &$[60;70)$& $[70;80)$&$[80;90)$\\
				\hline
				Giá trị đại diện & $45$ &$55$&$65$&$75$&$85$\\
				\hline
				Tần số & $5$&$8$&$25$&$20$&$2$\\
				\hline
			\end{tabular}
		\end{center}
		\begin{itemchoice}
			\itemch Số trung bình của mẫu số liệu trên là
			\[
			\overline{x}=\dfrac{1}{60}\left(5\cdot 45+8\cdot 55+25\cdot 65+20\cdot 75+2\cdot 85  \right)=66\,\,(\text{nghìn đồng}).
			\]
			Mệnh đề sai.
			\itemch Gọi $x_1;x_2;\ldots;x_{60}$ là số tiền mua sách xếp không giảm của $60$ khách hàng.\\
			Do $x_1;\ldots;x_5\in [40;50)$; $x_6;\ldots;x_{13}\in [50;60)$; $x_{14};\ldots;x_{38}\in [60;70)$; \\ $x_{39};\ldots;x_{58}\in [70;80)$; $x_{59};x_{60}\in [80;90)$.\\
			Suy ra trung vị của mẫu số liệu $x_1;\ldots; x_{60}$ là $\dfrac{1}{2}\left(x_{30}+x_{31}\right)\in [60;70)$.\\
			Ta xác định được $n=60$; $n_m=25$; $C=13$; $u_{m}=60$; $u_{m+1}=70$.\\
			Trung vị của mẫu số liệu ghép nhóm là
			\[
			M_\mathrm{e}= 60+\dfrac{\dfrac{60}{2}-13}{25}\left(70-60\right) =66{,}8\,\, (\text{nghìn đồng}).
			\]
			Mệnh đề đúng.
			\itemch Tứ phân vị thứ nhất của mẫu số liêụ xếp không giảm $x_1;x_2;\ldots;x_{60}$ là $x_{15}\in [60;70)$.\\
			Tứ phân vị thứ nhất $Q_1$ của mẫu số liệu ghép nhóm là
			\[
			Q_1 =60+ \dfrac{\dfrac{60}{4}-13}{25}\left(70-60\right) =60{,}8\,\, \text{(nghìn đồng)}.
			\]
			Mệnh đề đúng.
			\itemch Nhóm chứa mốt của mẫu số liệu ghép nhóm là $[60;70)$.\\
			Do đó $u_m=60$; $n_{m-1}=8; n_{m}=25; n_{m+1}=20$; $u_{m+1}-u_m=70-60=10$.\\
			Mốt của mẫu số liệu ghép nhóm là
			\[
			M_O=60+\dfrac{25-8}{(25-8)+(25-20)}\cdot\left(70-60\right)=\dfrac{745}{11}\approx 67{,}73 \,\, \text{(nghìn đồng)}.
			\]
			Mệnh đề sai.
		\end{itemchoice}
	}
\end{ex}

\begin{ex}%[1D5H2-3]
% nguon: Archive/2025-Sangtac/Khoi12/Data/1-KNTT/KNTT_1-OTGK1_De08.tex #5041

	Điểm kiểm tra giữa học kì 1 của lớp $11A1$ được thống kê như bảng dưới
	\begin{center}
		\begin{tabular}{|c|c|c|c|c|c|}
			\hline
			Khoảng điểm	& $[5;6)$ & $[6;7)$ & $[7;8)$ & $[8;9)$ & $[9;10]$ \\
			\hline
			Số học sinh	& $3$ & $5$ & $10$ & $6$ & $2$ \\
			\hline
		\end{tabular}
	\end{center}
	Xét tính đúng sai của các khẳng định sau
	\choiceTF
	{\True Mẫu số liệu đã cho là mẫu số liệu ghép nhóm}
	{\True Nhóm chứa mốt của số liệu là $[7;8)$}
	{Số học sinh đạt điểm khá (từ $8$ trở lên) là $18$}
	{\True Tứ phân vị thứ nhất của mẫu số liệu là $Q_1=6{,}7$}
	\loigiai{
		\begin{itemchoice}
			\itemch \textbf{Đúng.} Vì mẫu số liệu đã cho là mẫu số liệu ghép nhóm.
			\itemch \textbf{Đúng.} Vì  nhóm có tần số lớn nhất là $[7;8)$.
			\itemch \textbf{Sai.} Vì  số học sinh đạt điểm khá là $6+2=8$.
			\itemch \textbf{Đúng.} Vì gọi $x_1;\ x_2;\ldots;x_{26}$ là là mẫu số liệu được sắp xếp theo thứ tự không giảm.\\
			Trung vị của mẫu số liệu $\dfrac{x_{13}+x_{14}}{2}\in [7;8)$.\\
			Xét nửa mẫu số liệu bên trái $x_1;\ x_2;\ldots;x_{13}$ có trung vị $x_7\in[6;7)$.\\
			Khi đó $Q_1=6+\dfrac{\dfrac{26}{4}-3}{5}\cdot (7-6)=6{,7}$.
		\end{itemchoice}
	}
\end{ex}

\begin{bt}%[1D5V2-2]
% nguon: Archive/2025-Sangtac/Khoi12/Data/1-KNTT/KNTT_1-OTGK1_De08.tex #5047

	Bạn Chi vào website của một cửa hàng bán điện thoại tìm hiểu và đã thống kê số lượng một loại điện thoại theo giá bán cho kết quả như sau
	\begin{center}
		\begin{longtable}{|c|c|c|c|c|c|}
			\hline
			Giá tiền (triệu đồng) & $<2$ & $[2;4)$ & $[4;7)$ & $[7;13)$ & $[13;20]$\\
			\hline
			Số lượng & $20$ & $5$ & $11$ & $18$ & $21$\\
			\hline
		\end{longtable}
	\end{center}
	$50\%$ loại điện thoại có giá dưới bao nhiêu (làm tròn kết quả đến hàng phần trăm)? 
	
	\shortans{7{,}42}
	\loigiai{
		\begin{itemize}
			\item Dựa vào bảng, ta thấy có $20$ loại điện thoại mức giá dưới $2$ triệu đồng, $5$ loại điện thoại mức giá từ $2$ đến dưới $4$ triệu đồng, $11$ loại điện thoại mức giá từ $4$ đến dưới $7$ triệu đồng, $18$ loại điện thoại mức giá từ $7$ đến dưới $13$ triệu đồng, $21$ loại điện thoại mức giá từ $13$ đến $20$ triệu đồng.
			\item Ta có cỡ mẫu $n=20+5+11+18+21=75$.\\
			Gọi $x_1$, $x_2$, $\ldots$, $x_{75}$ là giá tiền của $75$ chiếc điện thoại và giả sử dãy này đã được sắp xếp theo thứ tự tăng dần.\\
			Khi đó, trung vị là $x_{38}$.\\
			Do giá trị $x_{38}$ thuộc nhóm $[7;13)$ nên nhóm này chứa trung vị.\\
			Do đó $a_4=7$, $m_4=18$, $m_1+m_2+m_3=36$; $a_5-a_4=5$ và ta có trung vị của mẫu số liệu trên là
			$$M_e=7+\dfrac{\dfrac{75}{2}-36}{18}\cdot 5\approx 7{,}42.$$
			Vậy $50\%$ điện thoại được bán có giá dưới $7{,}42$ triệu đồng.
		\end{itemize}
	}
\end{bt}

\begin{bt}%[1D5V1-2]
% nguon: Archive/2025-Sangtac/Khoi12/Data/1-KNTT/KNTT_1-OTGK1_De01.tex #5052

	Khi khảo sát chiều cao của $65$ em học sinh nam ở hai lớp 11A và 11B (đơn vị cm), ta thu được mẫu số liệu ghép nhóm được cho bảng sau:
	\begin{center}
		\begin{tabular}{|c|c|c|c|c|c|}
			\hline
			Nhóm & $[160; 163)$ & $[163; 166)$ &$[166; 169)$ &$[169; 172)$ & $[172; 174)$\\
			\hline
			Tần số & $x$ & 15 & 12 & 14 & $y$\\
			\hline
		\end{tabular}
	\end{center}
	Tính tổng $x+y$ biết tần số của nhóm $\left[ 172\,;174 \right)$ gấp $2$ lần tần số của nhóm $\left[ 160;,163 \right)$.
	
	\shortans{24}
	\loigiai{
		Ta có $$\heva{
			& x+15+12+14+y=65 \\
			& y=2x}\Leftrightarrow \heva{
			& x+y=24 \\
			& y=2x }\Leftrightarrow \heva{			& x=8 \\
			& y=16.}$$
	}
\end{bt}

\begin{ex}%[1D5H1-2]
% nguon: Archive/2025-Sangtac/Khoi12/Data/1-KNTT/KNTT_1-OTGK1_De07.tex #5057
 
	Cho mẫu số liệu ghép nhóm về tuổi thọ (đơn vị tính là năm) của một loại bóng đèn mới như
	sau
	\begin{center}		
		\begin{tabular}{|l|l|l|l|l|}
			\hline Tuổi thọ & {$[2 ; 3{,}5)$} & {$[3{,}5 ; 5)$} & {$[5 ; 6{,}5)$} & {$[6{,}5 ; 8)$} \\
			\hline Số bóng đèn & $8$ & $22$ & $35$ & $15$ \\
			\hline
		\end{tabular}
	\end{center}
	\choiceTF
	{\True Nhóm chứa mốt của mẫu số liệu là $[5; 6{,}5)$}
	{Độ dài lớn nhất của các nhóm ghép là $2$}
	{\True Có $80$ bóng đèn được thống kê}
	{\True Mốt của mẫu số liệu ghép nhóm là $5{,}59$}	
	\loigiai{
		\begin{itemchoice}
			\itemch Đúng.  Lớp chứa mốt là lớp có tần số lớn nhất, tức là lớp $[5; 6,5)$.
			\itemch Sai. Độ dài các lớp là $1{,}5.$ 
			\itemch Đúng. Ta có $8+22+35+15 =80$.
			\itemch Đúng. Ta có
			$L = 5$ (cận dưới của lớp $[5; 6{,}5)$). \\
			$f_1 = 35$  (tần số của lớp chứa mốt). \\
			$f_0 = 22$  (tần số của lớp ngay trước lớp chứa mốt  $[3{,}5; 5)$). \\
			$f_2 = 15$ (tần số của lớp ngay sau lớp chứa mốt $[6{,}5; 8)$). \\
			$h = 6,5 - 5 = 1,5$ (độ rộng của lớp chứa mốt).\\
			Ta có $M_0 = 5 + \dfrac{(35 - 22)}{(2 \cdot 35 - 22 - 15)} \cdot 1{,}5\approx 5 + 0,59 = 5{,}59$.\\			
			Vậy mốt của tuổi thọ bóng đèn là khoảng $5{,}59$ năm.
		\end{itemchoice}
	}
\end{ex}
